\documentclass[11pt]{amsart}
\usepackage[margin=1.1in]{geometry}
\usepackage{amsmath,amssymb,amsthm,mathtools}
\usepackage{booktabs,array}
\usepackage[T1]{fontenc}
\usepackage{lmodern}
\usepackage[hidelinks]{hyperref}
\usepackage{xcolor}

\newtheorem{theorem}{Theorem}[section]
\newtheorem{proposition}[theorem]{Proposition}
\newtheorem{lemma}[theorem]{Lemma}
\newtheorem{corollary}[theorem]{Corollary}
\theoremstyle{definition}
\newtheorem{definition}[theorem]{Definition}
\newtheorem{remark}[theorem]{Remark}
\numberwithin{equation}{section}
\raggedbottom
\emergencystretch=3em

\newcommand{\re}{\operatorname{Re}}
\newcommand{\im}{\operatorname{Im}}
\newcommand{\Sym}{\operatorname{Sym}^2}
\newcommand{\LS}{L_S}
\newcommand{\GR}{\Gamma_{\mathbb R}}
\newcommand{\GC}{\Gamma_{\mathbb C}}
\newcommand{\ZZ}{\mathbb Z}
\newcommand{\QQ}{\mathbb Q}
\newcommand{\RR}{\mathbb R}
\newcommand{\CC}{\mathbb C}
\newcommand{\FF}{\mathbb F}
\newcommand{\cO}{\mathcal O}
\newcommand{\repo}{\texttt{2026-11-number-theory-q-series/06-ramanujan-tau}}
\usepackage{url}\newcommand{\file}[1]{\path{#1}}
\newcommand{\GRHD}{\mathrm{GRH}(\Delta)}
\newcommand{\GRHS}{\mathrm{GRH}(\mathrm{Sym}^2\Delta)}

\begin{document}

\title[Ramanujan's tau function: signs at every $x$, excluded values]{Ramanujan's tau function: a sign theorem at every $x$\\ and excluded values}

\author{Raghav Sharma}
\email{raghavsharma1729@gmail.com}
\date{27 November 2026}
\subjclass[2020]{Primary 11F30; Secondary 11F66, 11M26, 11N05, 11D59, 11Y35, 11Y50}
\keywords{Ramanujan's tau function, Chebyshev's bias, Hecke eigenvalues, symmetric square $L$-function, explicit formula, verification of the Riemann hypothesis, Lehmer's conjecture, Thue equations, Lucas sequences}

\begin{abstract}
Let $\lambda(p)=\tau(p)p^{-11/2}$. Assuming only the Riemann hypothesis for $L(s,\Delta)$ (nothing about $\zeta$ or other $L$-functions, no linear independence), we prove that the Ces\`aro mean $\sum_{p\le x}\lambda(p)(1-p/x)$ is positive for every $x>5.5180546\ldots$ and nonpositive below, and that $\sum_{p\le t}\lambda(p)\log p\,(t-p)>0$ for every $t>4.2667919\ldots$, with an explicit lower bound of order $t^{3/2}$. Conversely, eventual nonnegativity of either sum implies the Riemann hypothesis for $L(s,\Delta)$. The positive sign comes from the prime squares, through $L(s,\mathrm{Sym}^2\Delta)$; to avoid assuming the Riemann hypothesis for it we prove, with a computer-assisted argument in ball arithmetic, that all its zeros in the critical strip with $|\mathrm{Im}\,s|\le 60$ are simple and lie on the critical line. Second, we prove that $x^{22}-3x^{11}y^2+y^4$ takes neither of the values $461$, $-599$ at integers, and that $\tau(p^2)\ne601$ for every prime $p$. With the reduction of Balakrishnan, Craig, Ono and Tsai this gives $\tau(n)\notin\{\pm461,\pm599,\pm601\}$ for all $n$; the cases $461$, $-599$ and $601$ are not settled in arXiv:2405.16723v1. The value $\tau(p^2)=463$ is excluded subject to one uncertified PARI/GP Thue computation.
\end{abstract}

\maketitle

\noindent\fbox{\parbox{\dimexpr\linewidth-2\fboxsep-2\fboxrule}{\small
\textbf{Disclosure.} This paper was produced by AI (Anthropic's Claude models), directed by the author. The proofs and computations were checked by independent AI referees, each instructed to find errors and working with its own code (Section~\ref{sec:review}). The paper has \emph{not yet been refereed by a human expert}. The author takes responsibility for the content and welcomes corrections at raghavsharma1729@gmail.com. Code, certificates and referee reports: \url{https://github.com/raghavs1729/Research-Math}, directory \repo.}}

\medskip
\tableofcontents

\section{Introduction}\label{sec:intro}

Let $\Delta(z)=q\prod_{n\ge1}(1-q^n)^{24}=\sum_{n\ge1}\tau(n)q^n$ ($q=e^{2\pi i z}$) and $\lambda(n)=\tau(n)n^{-11/2}$. By Deligne, $\lambda(p)=\alpha_p+\beta_p$ with $|\alpha_p|=|\beta_p|=1$, $\alpha_p\beta_p=1$. This paper has two independent parts: an analytic one on the sign of sums of $\lambda(p)$ over primes, valid at \emph{every} $x$ beyond an explicit threshold, and a Diophantine one on values that $\tau(n)$ does not take.

\subsection{Part I: the sign of prime sums of \texorpdfstring{$\lambda(p)$}{lambda(p)}}
Sarnak, in a letter to Mazur \cite{Sarnak}, observed that under the Riemann hypothesis for $L(s,\Delta)$ the normalised sum $(\log x/\sqrt x)\sum_{p\le x}\lambda(p)$ has a limiting logarithmic distribution with mean $1$: the sum is biased towards positive values. The source of the bias is the contribution of the prime squares, governed by the symmetric square $L$-function, which has no pole at $s=1$. Koyama and Kurokawa \cite{KK} showed, under the Deep Riemann Hypothesis, that $\sum_{p\le x}\tau(p)p^{-6}\sim\frac12\log\log x$ (in the restatement of Aoki--Koyama \cite{AK}, $=\frac12\log\log x+c+o(1)$); Sheth \cite[Corollary~B]{Sheth} proved $\sum_{p\le x}\tau(p)p^{-6}=\frac12\log\log x+c+o(1)$ under the Riemann hypothesis for $L(s,\Delta)$ alone, for $x$ outside a set of finite logarithmic measure, and Koyama--Sheth \cite{KS} note the analogue for all newforms. In particular positivity of $\sum_{p\le x}\tau(p)p^{-6}$ on a set of natural density $1$ under the Riemann hypothesis for $L(s,\Delta)$ follows from Sheth's theorem (with the elementary remark in Hayani \cite[(1.5)]{Hayani}); we do \emph{not} claim that statement. Here we prove sign statements valid at every $x$ beyond an explicit threshold, for natural smoothings of the unweighted sum. The method is that of the author's earlier work on the races modulo $4$ and $3$ \cite{WR}, adapted to $\tau$; the new ingredients are listed after Theorem~\ref{thm:4}.

Write $\GRHD$ for the statement that every zero of $L(s,\Delta)=\sum\lambda(n)n^{-s}$ with $0<\re s<1$ has real part $\frac12$, and $\GRHS$ for the same statement for $L(s,\Sym\Delta)$. Put
\[
\Theta_1(t)=\sum_{p\le t}\lambda(p)\log p\,(t-p),\qquad R(x)=\sum_{p\le x}\lambda(p)\Bigl(1-\frac px\Bigr)=\frac1x\int_0^x\sum_{p\le u}\lambda(p)\,du .
\]
Let $t^*=4.2667919563\ldots$ be the zero of $\Theta_1$ in $(3,5)$ and $x^*=5.5180546318\ldots$ the zero of $R$ in $(5,7)$ (closed forms in \eqref{eq:tstar}). The explicit functions $L_A\ge L_B$ are defined in \eqref{eq:LA}; $L_B(10^6)\ge0.0274$ and $\lim_{t\to\infty}L_B(t)\ge0.3692$.

\begin{theorem}[integrated $\theta$-sum]\label{thm:1}
\textup{(a)} Assume $\GRHD$. Then $\Theta_1(t)\ge L_B(t)\,t^{3/2}$ for every $t\ge1$, and $\Theta_1(t)>0$ for every $t>t^*$; also $\Theta_1(t)\le0$ for $t\le t^*$.
\textup{(b)} Assume $\GRHD$ and $\GRHS$. Then moreover $\Theta_1(t)\ge L_A(t)\,t^{3/2}$ for every $t\ge1$, and $\lim L_A\ge0.5539$.
\end{theorem}

\begin{theorem}[Ces\`aro mean]\label{thm:2}
Assume $\GRHD$. Then $R(x)>0$ for every $x>x^*$, and $R(x)\le0$ for $x\le x^*$. More precisely, $R(x)\ge \Theta_1(x)/(x\log x)+0.6719\ge L_B(x)\sqrt x/\log x+0.6719$ for every $x\ge97$, and $R(x)>0.15\sqrt x/\log x$ for every $x\ge11$. If $\GRHS$ holds as well, then $R(x)>\sqrt x/(3\log x)$ for every $x\ge100$.
\end{theorem}

\begin{theorem}[equivalence; unconditional]\label{thm:3}
The following are equivalent: \textup{(a)} $\GRHD$; \textup{(b)} $R(x)>0$ for every $x>x^*$; \textup{(c)} $R(x)\ge0$ for all sufficiently large $x$; \textup{(d)} $\Theta_1(t)>0$ for every $t>t^*$; \textup{(e)} $\Theta_1(t)\ge0$ for all sufficiently large $t$.
\end{theorem}

\begin{theorem}[computer-assisted; no hypothesis]\label{thm:4}
$L(s,\Sym\Delta)$ has exactly $41$ zeros with $0<\im s\le60$ and $0<\re s<1$; they are simple and have real part $\frac12$ (the first is $\frac12+i\,7.0702241402\ldots$), and $L(s,\Sym\Delta)$ has no zero on the real segment $(0,1)$.
\end{theorem}

Theorem~\ref{thm:4} is what removes $\GRHS$ from Theorems~\ref{thm:1}(a), \ref{thm:2} and from the implications (a)$\Rightarrow$(b), (a)$\Rightarrow$(d) of Theorem~\ref{thm:3}: without it these hold under ``$\GRHD$ and the Riemann hypothesis for $L(s,\Sym\Delta)$ up to height $60$''. The implications (b),(c),(d),(e)$\Rightarrow$(a) and Theorem~\ref{thm:1}(b) do not use Theorem~\ref{thm:4}.

Deligne's bound alone does not suffice for the prime-square term: it gives only $|\cdot|\le\frac43t^{3/2}$, against a bias of $\frac23t^{3/2}$. The new ingredients relative to \cite{WR} are: Lemma~\ref{lem:decomp}, the treatment of the symmetric square without assuming its Riemann hypothesis (Lemmas~\ref{lem:lowzeros}, \ref{lem:highzeros} and Theorem~\ref{thm:4}), the closed form for $\sigma_0^2$ certified from a theta integral, and the handling of the initial negative patch in Section~\ref{sec:proof2}. The point of Lemma~\ref{lem:decomp} is that all even prime powers are carried by $\psi_S-\psi$, so that the main term comes from $\psi$, and the symmetric square needs only a partial Riemann hypothesis. The sizes involved are: bias $\frac23$; zeros of $L(s,\Delta)$ at most $\sigma_0^2:=2(\Lambda'/\Lambda)(1)=0.1122521\ldots$ (which equals $\sum_\rho|\rho|^{-2}$ under $\GRHD$), a closed form certified without a list of zeros; zeros of $L(s,\Sym\Delta)$ above height $60$ at most $0.1846$ (an unconditional bound, Lemma~\ref{lem:highzeros}).

The completeness test in the proof of Theorem~\ref{thm:4} (Lemma~\ref{lem:V3}) is a variant of the method of Hiary, Ireland and Kyi \cite{HIK}, who evaluate the positive Hadamard sum at a real point from the Euler product and compare it with the contribution of the zeros found on the line; their examples include $L(s,\Delta)$ but no symmetric square. Lemma~\ref{lem:V3} was written before that paper was read, and we do not claim it as a new method; what we claim is Theorem~\ref{thm:4} for $L(s,\Sym\Delta)$. Yoshida \cite{Yoshida88} computed zeros of $L(s,\Sym\Delta)$ (as reported in \cite{Yoshida95}, where he describes his method as lacking rigorous error control; \cite{Yoshida88} was not seen by us). We found no rigorous verification of the Riemann hypothesis for $L(s,\Sym\Delta)$ to any height in the literature we searched; but the LMFDB could not be queried (access was blocked), and Booker's work on rigorous computation of general $L$-functions \cite{Booker} and the announced work of Bober, Booker, Costa, Lee, Platt and Sutherland on motivic $L$-functions were not read.

\emph{Relation to the literature.} We found no statement for $\tau$ valid for every $x$ beyond an explicit $x_0$, and no equivalence of the kind in Theorem~\ref{thm:3} for a sum over primes of the coefficients of a cusp form. The nearest results are: Johnston's all-$x$ equivalence for $\zeta$, RH $\iff\int_2^x(\pi(t)-\mathrm{li}(t))\,dt<0$ for all $x>2$ \cite{Johnston} (quoted from Suzuki's introduction and the arXiv abstract); Hayani's Theorem~1.4 \cite{Hayani}, a Ces\`aro mean at every $x$ under GRH for Dirichlet characters, with a non-explicit constant; Suzuki's equivalences \cite{Suzuki1} for $\zeta$ and Dirichlet $L$-functions; and Suzuki's equivalence for the extended Selberg class \cite{Suzuki2}, which covers $L(s,\Delta)$ in principle but concerns a screw function (a prime-power sum plus archimedean terms) and does not treat $\tau$ explicitly. Devin \cite{Devin} gives the general framework for the limiting distribution. ``Not found'' refers to the searches recorded in \file{tau-race/LITERATURE.md} of the repository.

\subsection{Part II: values not taken by \texorpdfstring{$\tau$}{tau}}
Lehmer asked whether $\tau(n)\ne0$ for all $n$. Balakrishnan, Craig and Ono \cite{BCO}, and with Tsai \cite{BCOT}, studied which odd values $\pm\ell$ can occur. By \cite{BCO}, $\tau(n)\notin\{\pm1,\pm3,\pm5,\pm7,\pm691\}$ for $n>1$ (as stated in its abstract); \cite{BCOT} reduced the equation $\tau(n)=\pm\ell$, for an odd prime $\ell\nmid\tau(\ell)$, to $n=p^{d-1}$ with $p$ and $d$ odd primes, $d\mid\ell(\ell^2-1)$. Goto \cite[Theorem~0.3]{Goto} states: for an odd prime $\ell<1000$ and $\varepsilon=\pm1$, if $\tau(n)=\varepsilon\ell$ then $n=p^4$ with $\varepsilon\ell\in\{461,-599\}$. His Section~2.3 says: ``When $d=5$, the equation (2.1) for $a=281,461,-599,-919$ cannot be ruled out by the congruences''; he then rules out $281$ and $-919$ (by a method of Dembner and Jain \cite{DJ}) and leaves $461$ and $-599$ open. Odd values of $\tau$ were also studied by Bennett, Gherga, Patel and Siksek \cite{BGPS}, Dembner and Jain \cite{DJ}, and Hanada and Madhukara \cite{HM}.

\begin{theorem}\label{thm:p4}
Neither $x^{22}-3x^{11}y^2+y^4=461$ nor $x^{22}-3x^{11}y^2+y^4=-599$ has a solution in integers. Consequently $\tau(p^4)\notin\{461,-599\}$ for every prime $p$.
\end{theorem}

\begin{theorem}\label{thm:p2}
$\tau(p^2)\ne601$ for every prime $p$. Moreover $\tau(p^2)\ne463$ for every prime $p$, \emph{provided} the PARI/GP computation described in Section~\ref{sec:463} (one Thue equation of degree $11$ with right-hand side $\pm1$, solved with an uncertified class group and unit computation) is correct.
\end{theorem}

\begin{corollary}\label{cor:values}
$\tau(n)\notin\{\pm461,\pm599,\pm601\}$ for every $n\ge1$, and $\tau(n)\ne-463$. If the computation of Section~\ref{sec:463} is correct, also $\tau(n)\ne463$.
\end{corollary}

Corollary~\ref{cor:values} uses the literature (the reduction of \cite{BCOT} and classical congruences); it uses Goto's paper only as an exposition of a standard step (\cite[\S3.1]{Goto}). The values $463$ and $601$ appear because of the case $d=3$ in the arXiv version of Goto's argument: there the proof of Theorem~0.3 lists no pair with $d=3$; the pair $(601,3)$ passes every congruence of his Lemma~A.1 as printed, so its omission is an error in that list; separately, the mod-$12$ line ``$\tau(p^{d-1})\equiv0,1,d\pmod{12}$'' of Lemma~A.1 is false for $d=3$ ($\tau(5^2)=-25499225\equiv7\pmod{12}$; the true set of residues is $\{1,3,7,9\}$) and, according to one of our referees, also for $d=11$, and it wrongly excludes $463\equiv7\pmod{12}$; the pair $(-827,3)$ is not excluded by any of our congruences either (Section~\ref{sec:goto}). The published version of Goto's paper \cite{Goto} (Arch.\ Math.\ 125 (2025)) was not accessible to us; its abstract refers to an explicit finite set of odd primes below $1000$, which we have not seen. We therefore claim only that the cases treated here are not settled by arXiv:2405.16723v1.

For $461$ the condition that $V=\tau(p)^2$ is a square is essential, and it is what is new in the argument relative to the treatments we have read: these (\cite{BCOT,DJ,Goto}, \cite{Angdinata}) pass from $\tau(p^4)=U^2-3UV+V^2$, $(U,V)=(p^{11},\tau(p)^2)$, to the hyperelliptic equation $u^2=5x^{22}+4\alpha$, which forgets that $V$ is a square. For $-599$ the eleventh-power condition already suffices, by the method Goto uses for $281$ and $-919$ (Remark~\ref{rem:which}). As a by-product we settle four cases left open by Hanada and Madhukara for weights $16$ and $26$ (Proposition~\ref{prop:HM}).

\subsection{Hypotheses and conventions}
Results conditional on a Riemann hypothesis say so in their statements; no statement of this paper uses a hypothesis on $\zeta$, a linear-independence hypothesis or the Deep Riemann Hypothesis. A Lean proof of GRH for Dirichlet $L$-functions was announced on 2026-11-20; it is unrefereed, does not concern $L(s,\Delta)$, and is not used here. ``Certified'' means decided in ball arithmetic (Arb \cite{Arb}, through python-flint 0.9.0) or with exact integers, every comparison rigorous, correct given FLINT/Arb and the named program; in Part~II it also covers results proved by PARI's \texttt{bnfcertify} or by \texttt{thueinit} with certification flag, correct given PARI \cite{PARI}. ``Numerical'' means floating point. File names refer to the directory \repo{} of the repository \cite{Repo}.

\part*{Part I. A sign theorem at every \texorpdfstring{$x$}{x}}
\section{Notation and quoted inputs}\label{sec:notation}

Throughout, $p$ runs over primes, $\lambda(p)=\alpha_p+\beta_p$ and $c_k(p)=\alpha_p^k+\beta_p^k\in[-2,2]$, so that $c_1=\lambda(p)$, $c_2=\lambda(p)^2-2$, $c_{k+1}=\lambda(p)c_k-c_{k-1}$. Let
\begin{align*}
L(s,\Delta)&=\sum_{n\ge1}\lambda(n)n^{-s}=\prod_p(1-\alpha_pp^{-s})^{-1}(1-\beta_pp^{-s})^{-1},\qquad \Lambda(s)=(2\pi)^{-s}\Gamma(s+\tfrac{11}2)L(s,\Delta),\\
L_S(s)&=L(s,\Sym\Delta)=\prod_p\bigl[(1-\alpha_p^2p^{-s})(1-p^{-s})(1-\beta_p^2p^{-s})\bigr]^{-1}=\sum_{n\ge1}a_S(n)n^{-s},\\
\gamma_S(s)&=\GR(s+1)\GC(s+11)=\pi^{-(s+1)/2}\Gamma\bigl(\tfrac{s+1}2\bigr)\cdot2(2\pi)^{-(s+11)}\Gamma(s+11),\qquad \Lambda_S=\gamma_SL_S .
\end{align*}
We write $\Lambda_\Delta(p^k)=c_k(p)\log p$, $\Lambda_{\rm sym}(p^k)=(c_{2k}(p)+1)\log p$, $\Lambda$ for von Mangoldt's function, $\psi_\Delta,\psi_S,\psi$ for the corresponding summatory functions and $\theta(y)=\sum_{p\le y}\log p$. Further
\[
\begin{gathered}\Theta(t)=\sum_{p\le t}\lambda(p)\log p,\qquad \Theta_1(t)=\int_0^t\Theta=\sum_{p\le t}\lambda(p)\log p\,(t-p),\\ D_0(x)=\sum_{p\le x}\lambda(p),\qquad D_1(x)=\sum_{p\le x}\lambda(p)p,\end{gathered}
\]
so that $R(x)=D_0(x)-D_1(x)/x$. $\mathrm{RH}_S(H)$ is the statement that every zero $\rho$ of $L_S$ with $0<\re\rho<1$ and $|\im\rho|\le H$ has real part $\frac12$.

\medskip\noindent\textbf{Quoted inputs.}
\begin{itemize}
\item[(F1)] (Deligne \cite{Deligne}.) $|\alpha_p|=1$.
\item[(F2)] (Hecke.) $\Lambda$ is entire, $\Lambda(s)=\Lambda(1-s)$, bounded in vertical strips.
\item[(F3)] (Rankin \cite{Rankin}, Shimura \cite{Shimura}, Gelbart--Jacquet \cite{GJ}.) $\Lambda_S$ is entire, $\Lambda_S(s)=\Lambda_S(1-s)$, and $L_S$ has finite order. We read this in two sources. Zagier \cite[eq.~(4), p.~109 and Corollary~1, p.~111]{Zagier}: for a normalised eigenform $f$ of weight $k$ on $\mathrm{SL}_2(\ZZ)$ and $D_f(s)=\prod_p(1-\alpha_p^2p^{-s})^{-1}(1-\alpha_p\bar\alpha_pp^{-s})^{-1}(1-\bar\alpha_p^2p^{-s})^{-1}$ (here $\alpha_p\bar\alpha_p=p^{k-1}$), $D_f^*(s)=2^{-s}\pi^{-3s/2}\Gamma(s)\Gamma(\frac{s-k+2}2)D_f(s)=D_f^*(2k-1-s)$ and $D_f$ is entire. With $s=s'+k-1$ one has $D_f(s)=L_S(s')$ and $D_f^*(s'+k-1)=\frac12\pi^{-(k-2)/2}\GR(s'+1)\GC(s'+k-1)L_S(s')$, a constant multiple of $\Lambda_S(s')$. Balkanova--Frolenkov \cite[(2-4)--(2-5)]{BF} state the same with $L_\infty(s)=\pi^{-3s/2}\Gamma(\frac{s+1}2)\Gamma(\frac{s-1}2+k)\Gamma(\frac s2+k)$ (weight $2k$), a constant multiple of $\gamma_S$ by the duplication formula. The originals of Rankin, Shimura and Gelbart--Jacquet were not consulted; the finite order of $L_S$ is cited, not checked. Numerical confirmation: PARI's \texttt{lfuncheckfeq} gives $-67$ bits (and $-111$ bits in an independent set-up; $-2$ bits with a wrong gamma shift), and the identity $\Phi(1/y)=y\Phi(y)$ of Lemma~\ref{lem:V0} was checked directly to $10^{-83}$ at three points in ball arithmetic.
\item[(F4)] (Hadamard; \cite[Theorem~5.6, Proposition~5.7]{IK}.) For $\mathcal L\in\{L(\cdot,\Delta),L_S\}$ with completed function $\Lambda_{\mathcal L}$: $\re(\Lambda_{\mathcal L}'/\Lambda_{\mathcal L})(s)=\sum_\rho\re\frac1{s-\rho}$, absolutely convergent, and $N(T+1)-N(T)\ll\log(|T|+2)$. Neither function vanishes at $s=0,1$ ($L(1,\Delta)=0.8393\ldots$ and $L_S(1)=0.6317\ldots$ are certified).
\item[(F5)] \cite[Proposition~5.7(2)]{IK}: for $-\frac12\le\sigma\le2$, $(\mathcal L'/\mathcal L)(s)=\sum_{|s-\rho|<1}(s-\rho)^{-1}+O(\log(|s|+2))$; to the left of the strip, at distance $\ge\frac14$ from the trivial zeros, $(\mathcal L'/\mathcal L)(s)\ll\log(|s|+2)$ by the functional equation.
\item[(F6)] $|\theta(y)-y|<0.2\,y/\log^2y$ for $y\ge3\,594\,641$: Broadbent--Kadiri--Lumley--Ng--Wilk \cite[arXiv:2002.11068v2, Corollary~11.2, eq.~(3.39), row $M_2=0.2$]{BKLNW} (the journal numbering may differ); the same inequality is \cite[Theorem~5.2]{Dusart}, an unrefereed preprint relying on unpublished zero verifications. BKLNW note that the proof of the corresponding theorem in Dusart's published paper is incorrect and re-prove it. For very large $y$ BKLNW use explicit zero-density inputs which Fiori, Kadiri and Swidinsky \cite{FKS} call unreliable; we found no erratum. This input is not needed for the sign statements: see Remark~\ref{rem:robust}, which uses only $\theta(y)\le1.01624\,y$ and $\theta(y)\ge y(1-1/\log y)$ for $y\ge10^7$, consequences of Rosser--Schoenfeld \cite{RS62} (as quoted in \cite[Tables~1--2]{BKLNW}; original not seen).
\end{itemize}

\medskip\noindent\textbf{Certified constants} (\file{tau-race/data/constants.json}, \file{tau-race/data/bounds.txt}):
\begin{align*}
\sigma_0^2&=2(\Lambda'/\Lambda)(1)=0.112252140471908526745583426253\ldots\\&\qquad(=\textstyle\sum_\rho|\rho|^{-2}\text{ under }\GRHD),\\
a_0&=(L'/L)(0,\Delta)=0.170657847591640\ldots,\\ a_{-1}&=(L'/L)(-1,\Delta)\in[0.281796,0.281803],\\
(\Lambda_S'/\Lambda_S)(2)&\in[0.2941460,0.2941565],\\ C_S&:=\tfrac{12}5(\Lambda_S'/\Lambda_S)(2)\le0.70598,\\ K_S&\le0.34176,\qquad T_{60}\le0.184605,\\
\varepsilon_1&=0.2/\log^2(10^7)=7.69844\ldots\cdot10^{-4},\qquad c_\theta=1+\varepsilon_1,
\end{align*}
where $K_S$ is an upper bound for $(L_S'/L_S)(-2)$ and $T_{60}$ is the bound of Lemma~\ref{lem:highzeros}. The thresholds are
\begin{gather}\label{eq:tstar}
t^*=\frac{2\lambda(2)\log2+3\lambda(3)\log3}{\lambda(2)\log2+\lambda(3)\log3}=4.266791956347541\ldots,\\ x^*=\frac{D_1(5)}{D_0(5)}=5.518054631888274\ldots .
\end{gather}
For $t\ge1$ define
\begin{equation}\label{eq:LA}
\begin{aligned}
E_{\rm odd}(t)&=\tfrac34t^{-1/6}+\tfrac56t^{-3/10}+t^{-5/14}\max\bigl(0,\tfrac12(\log t/\log2-5)\bigr),\\
L_A(t)&=\tfrac23\bigl(1-\varepsilon_1\mathbf 1_{t>10^{14}}\bigr)-\sigma_0^2-a_0t^{-1/2}-\bigl(\tfrac45\sqrt2+C_S\bigr)t^{-1/4}\\&\qquad-2t^{-1}-K_St^{-3/2}-2c_\theta E_{\rm odd}(t),\\
L_B(t)&=L_A(t)-T_{60}.
\end{aligned}
\end{equation}
$L_A,L_B$ are nondecreasing on $[527,10^{14}]$ and on $(10^{14},\infty)$, and (certified) $L_A(10^6)\ge0.2120$, $L_A(10^7)\ge0.3483$, $L_A(10^{14}+)\ge0.5458$ (here $f(10^{14}+)$ denotes the right limit), $\lim L_A\ge0.5539$; $L_B(10^6)\ge0.0274$, $L_B(10^7)\ge0.1637$, $L_B(10^{14}+)\ge0.3612$, $\lim L_B\ge0.3692$.

\section{The decomposition and the explicit formulas}\label{sec:explicit}

\begin{lemma}[decomposition; uses only (F1)]\label{lem:decomp}
For $u>0$,
\[
\Theta(u)=\psi_\Delta(u)-\bigl[\psi_S(\sqrt u)-\psi(\sqrt u)\bigr]-\sum_{k\ge3\ \mathrm{odd}}\Theta^{(k)}(u^{1/k}),\qquad \Theta^{(k)}(y):=\sum_{p\le y}c_k(p)\log p,
\]
and $|\Theta^{(k)}(y)|\le2\theta(y)$. Consequently, for $t>0$,
\begin{equation}\label{eq:decomp}
\Theta_1(t)=\psi_{1,\Delta}(t)-J_S(t)+J_\zeta(t)-\sum_{k\ge3\ \mathrm{odd}}\int_0^t\Theta^{(k)}(u^{1/k})\,du,
\end{equation}
where $\psi_{1,\Delta}(t)=\sum_n\Lambda_\Delta(n)(t-n)_+$, $J_S(t)=\sum_n\Lambda_{\rm sym}(n)(t-n^2)_+$ and $J_\zeta(t)=\sum_n\Lambda(n)(t-n^2)_+=\int_0^{\sqrt t}\psi(y)\,2y\,dy$.
\end{lemma}

\begin{proof}
$\psi_\Delta(u)=\sum_{k\ge1}\sum_{p^k\le u}c_k(p)\log p$. The terms with $k=2j$ sum to $\sum_{p^j\le\sqrt u}c_{2j}(p)\log p=\sum_{p^j\le\sqrt u}(\Lambda_{\rm sym}(p^j)-\log p)=\psi_S(\sqrt u)-\psi(\sqrt u)$, because the local factor of $L_S$ at $p$ has roots $\alpha_p^2,1,\beta_p^2$. The term $k=1$ is $\Theta(u)$, and $|c_k|\le2$ by (F1). Integrate over $[0,t]$; $\int_0^t\psi_S(\sqrt u)\,du=\sum_n\Lambda_{\rm sym}(n)(t-n^2)_+$, and similarly for $\psi$.
\end{proof}

Thus all even prime powers are carried by $\psi_S$ and $\psi$: no sign-indefinite even power has to be bounded trivially, and the main term $+\frac23t^{3/2}$ comes from $J_\zeta$, i.e.\ from $\psi$.

\begin{lemma}[moving the contour]\label{lem:contour}
Let $\mathcal L$ be $L(\cdot,\Delta)$ or $L_S$, $m\in\{1,2\}$, $t>1$ and $c>1/m$. Then $\frac1{2\pi i}\int_{(c)}(-\mathcal L'/\mathcal L)(ms)\,\frac{t^{s+1}}{s(s+1)}\,ds$ equals the sum of the residues of the integrand in $\re s<c$: at $s=\rho/m$ ($\rho$ a nontrivial zero), at $s=w_0/m$ ($w_0$ a trivial zero), at $s=0$ and at $s=-1$. The sum over $\rho$ converges absolutely.
\end{lemma}

\begin{proof}
(i) By (F4), (F5) there are $T_k\in[k,k+1]$ with $(\mathcal L'/\mathcal L)(\sigma\pm iT_k)\ll\log^2T_k$ uniformly for $-\frac12\le\sigma\le2$: there are $\ll\log k$ ordinates in $[k,k+1]$, so some $T_k$ has distance $\gg1/\log k$ from all of them. (ii) For $\re w\le-\frac12$ at distance $\ge\frac14$ from the trivial zeros, $(\mathcal L'/\mathcal L)(w)=-(\mathcal L'/\mathcal L)(1-w)-(\gamma'/\gamma)(w)-(\gamma'/\gamma)(1-w)\ll\log(|w|+2)$, since the first term is bounded and $\gamma'/\gamma$ is a constant plus digamma functions, $\digamma(z)\ll\log(|z|+2)$ at distance $\ge\frac18$ from its poles. Integrate over the rectangle with vertices $c\pm iT_k/m$, $-U'\pm iT_k/m$, where $mU'\in\mathbb N$ for $L(\cdot,\Delta)$ (trivial zeros at half-integers) and $mU'\in\mathbb N+\frac12$ for $L_S$ (trivial zeros at integers). On the horizontal sides $|s(s+1)|\ge T_k^2/m^2$, so these integrals are $\ll_t\log^2T_k/T_k^2$ uniformly in $U'$ (as $t>1$). On the left side, by (ii), the integral is $\ll t^{1-U'}\log(U'+2)/U'$. Let $U'\to\infty$, then $k\to\infty$. The residue at $s=\rho/m$ has modulus $\le m\,\mu_\rho t^2/|\rho|^2$ ($\mu_\rho$ the multiplicity), and $\sum_\rho|\rho|^{-2}<\infty$.
\end{proof}

\begin{lemma}[explicit formula for $\psi_{1,\Delta}$]\label{lem:psi1}
For $t\ge1$,
\[
\psi_{1,\Delta}(t)=-\sum_\rho\frac{t^{\rho+1}}{\rho(\rho+1)}-a_0t+a_{-1}-\sum_{n\ge0}\frac{t^{-9/2-n}}{(n+\frac{11}2)(n+\frac92)},
\]
the sum over the nontrivial zeros of $L(s,\Delta)$ (with multiplicity) being absolutely convergent. Under $\GRHD$, $\sum_\rho|\rho(\rho+1)|^{-1}<\sum_\rho|\rho|^{-2}=\sigma_0^2$, and
\begin{equation}\label{eq:psi1bound}
\psi_{1,\Delta}(t)\ge-\sigma_0^2t^{3/2}-a_0t\qquad(t\ge1).
\end{equation}
\end{lemma}

\begin{proof}
For $t>1$, $\psi_{1,\Delta}(t)=\frac1{2\pi i}\int_{(2)}(-L'/L)(s,\Delta)\frac{t^{s+1}}{s(s+1)}ds$. Apply Lemma~\ref{lem:contour} with $m=1$; the trivial zeros are the simple poles of $\Gamma(s+\frac{11}2)$ at $s=-\frac{11}2-n$. The residues are: $-t^{\rho+1}/(\rho(\rho+1))$ at $\rho$; $-a_0t$ at $s=0$ ($L(0,\Delta)=\Lambda(1)/\Gamma(\frac{11}2)\ne0$); $a_{-1}$ at $s=-1$ ($L(-1,\Delta)\ne0$); and $-t^{-9/2-n}/((n+\frac{11}2)(n+\frac92))$ at $-\frac{11}2-n$. Both sides are continuous at $t=1$. By (F4) at $s=1$, $(\Lambda'/\Lambda)(1)=\sum_\rho\re(1-\rho)^{-1}$, and under GRH $\re(1-\rho)^{-1}=\frac12|\rho|^{-2}$; also $|\rho+1|>|\rho|$. Finally $\sum_{n\ge0}((n+\frac{11}2)(n+\frac92))^{-1}=\frac29$ (telescoping), so for $t\ge1$ the last two terms are $\ge a_{-1}-\frac29>0.059$.
\end{proof}

(As checks, with PARI at 38 digits: the formula at $t=1$ holds to $4\cdot10^{-38}$, but by the functional equation this identity reduces to $\digamma(\frac{11}2)-\digamma(\frac92)=\frac29$ and only tests the bookkeeping of trivial zeros; the meaningful test is a direct comparison with $277$ pairs of zeros at $t=2.5,5.5,10$, with differences $\le6.6\cdot10^{-4}$, inside the tail bound $0.0103\,t^{3/2}$.)

\begin{lemma}[explicit formula for $J_S$]\label{lem:JS}
For $t\ge1$,
\[
J_S(t)=-\sum_\rho\frac{2t^{1+\rho/2}}{\rho(\rho+2)}-a_{S,0}t+2t^{1/2}+(L_S'/L_S)(-2)-\sum_{w}\frac{m_w}2\,\frac{t^{1-w/2}}{(w/2)(w/2-1)},
\]
where $\rho$ runs over the nontrivial zeros of $L_S$, $a_{S,0}=(L_S'/L_S)(0)$, and $w$ runs over the integers $w\ge3$ such that $-w$ is a trivial zero of $L_S$, with multiplicity $m_w$ ($m_w=1$ for $w\in\{3,5,7,9\}$ and for even $w\ge12$; $m_w=2$ for odd $w\ge11$). All terms of the last sum are positive. Consequently, with $Z_S(t):=\bigl|\sum_\rho2t^{1+\rho/2}/(\rho(\rho+2))\bigr|$,
\begin{equation}\label{eq:JSbound}
J_S(t)\le Z_S(t)+2t^{1/2}+K_S\qquad(t\ge1),\ \text{provided }\mathrm{RH}_S(2)\text{ holds}.
\end{equation}
\end{lemma}

\begin{proof}
For $t>1$ and $c>\frac12$,
\[
J_S(t)=\frac1{2\pi i}\int_{(c)}(-L_S'/L_S)(2s)\frac{t^{s+1}}{s(s+1)}\,ds
\]
(termwise, $\frac1{2\pi i}\int(t/n^2)^s\frac{t\,ds}{s(s+1)}=t(1-n^2/t)_+$). Apply Lemma~\ref{lem:contour} with $m=2$. The trivial zeros of $L_S$ are the poles of $\gamma_S$: $\Gamma(\frac{s+1}2)$ gives $-1,-3,-5,\dots$ and $\Gamma(s+11)$ gives $-11,-12,-13,\dots$. Residues: $-2t^{1+\rho/2}/(\rho(\rho+2))$ at $s=\rho/2$; $-a_{S,0}t$ at $s=0$; $2t^{1/2}$ at $s=-\frac12$ (the simple trivial zero $-1$); $(L_S'/L_S)(-2)$ at $s=-1$ (regular: $L_S(-2)\ne0$); and $-\frac{m_w}2t^{1-w/2}/((w/2)(w/2-1))$ at $s=-w/2$, $w\ge3$. For \eqref{eq:JSbound}: the trivial-zero sum is subtracted and positive, and $(L_S'/L_S)(-2)=-(\gamma_S'/\gamma_S)(-2)-(\gamma_S'/\gamma_S)(3)-(L_S'/L_S)(3)$ with $|(L_S'/L_S)(3)|\le-3(\zeta'/\zeta)(3)$, which gives $K_S\le0.34176$. It remains to see $a_{S,0}\ge0$. By the functional equation $a_{S,0}=-(\gamma_S'/\gamma_S)(0)-(\Lambda_S'/\Lambda_S)(1)$ with $-(\gamma_S'/\gamma_S)(0)=1.04024\ldots$, and $(\Lambda_S'/\Lambda_S)(1)=\sum_\rho\re(1-\rho)^{-1}\le3\sum_\rho\re(2-\rho)^{-1}=3(\Lambda_S'/\Lambda_S)(2)\le0.88247$: for a zero on the line the ratio of the two terms is $(\frac94+\gamma^2)/(3(\frac14+\gamma^2))\le3$, and for a zero $\beta+i\gamma$ with $|\gamma|>2$ the function $a\mapsto a/(a^2+\gamma^2)$ increases on $(0,|\gamma|)$, so $\re(1-\rho)^{-1}\le\re(2-\rho)^{-1}$; $\mathrm{RH}_S(2)$ covers the remaining zeros. Hence $a_{S,0}\ge0.1577>0$ (certified).
\end{proof}

(The identity at $t=1$ again reduces to a digamma identity and tests only the list of trivial zeros; the true values are $a_{S,0}=0.94103$ and $(L_S'/L_S)(-2)=-0.22825$. Direct comparisons with truncated zero sums are recorded in \file{tau-race/AUDIT-NUMERICS.md}, item~6.)

\section{Zeros of the symmetric square, and inputs on primes}\label{sec:sym2zeros}

\begin{lemma}[low zeros]\label{lem:lowzeros}
If $\mathrm{RH}_S(H)$ holds for some $H\ge2$, then $\sum_{|\im\rho|\le H}2/|\rho(\rho+2)|\le C_S$; under $\GRHS$ the same bound holds for the sum over all zeros.
\end{lemma}

\begin{proof}
For $\rho=\frac12+i\gamma$, with $g=\gamma^2$: $2/|\rho(\rho+2)|=2((\frac14+g)(\frac{25}4+g))^{-1/2}$ and $\re(2-\rho)^{-1}=\frac32(\frac94+g)^{-1}$. The quotient $\frac43(\frac94+g)((\frac14+g)(\frac{25}4+g))^{-1/2}$ is at most $\frac{12}5$, its value at $g=0$, since $25(g+\frac94)^2\le81(g+\frac14)(g+\frac{25}4)\iff0\le56g^2+414g$. All terms of $\sum_\rho\re(2-\rho)^{-1}=(\Lambda_S'/\Lambda_S)(2)$ are positive.
\end{proof}

$(\Lambda_S'/\Lambda_S)(2)=(\gamma_S'/\gamma_S)(2)+(L_S'/L_S)(2)$ is computed from the Euler product $-(L_S'/L_S)(s)=\sum_p\log p\,[(c_2x-2x^2)/(1-c_2x+x^2)+x/(1-x)]$, $x=p^{-s}$, $c_2=\tau(p)^2/p^{11}-2$ (exact rationals for $p\le10^7$), with tail $\le3\sum_{n>P}\log n/(n^2-1)\le3.03(\log P+1)/P$ at $\re s=2$: the bracket is $\sum_{k\ge1}(\alpha_p^{2k}+1+\beta_p^{2k})x^k$, of modulus $\le3/(p^2-1)$; $1/(n^2-1)\le1.01/n^2$ for $n\ge11$; and $\sum_{n>P}\log n/n^2\le\int_P^\infty\log u\,u^{-2}du=(\log P+1)/P$.

\begin{lemma}[high zeros; unconditional]\label{lem:highzeros}
Let $H>H'>0$, $c=\sqrt2/\arctan((H-H')/\sqrt2)$ and $(x_1,x_2,x_3)=(\frac32,\frac{13}2,7)$. For the zeros $\rho=\beta+i\gamma$ of $L_S$ in the critical strip,
\[
\sum_{|\gamma|>H}\frac2{\gamma^2}\le T(H,H'):=2c\Bigl[\frac32\,\frac{\log(H'/2\pi)+1}{H'}+\frac{\sum_jx_j^2}{3H'^3}+\frac{\sum_j(24x_j)^{-1}}{2H'^2}+\frac{6\log\zeta(2)}{H'^2}\Bigr].
\]
In particular $T_{60}:=T(60,51.45)\le0.184605$, and $\bigl|\sum_{|\gamma|>60}2t^{1+\rho/2}/(\rho(\rho+2))\bigr|\le T_{60}t^{3/2}$ for $t\ge1$.
\end{lemma}

\begin{proof}
Put $F(T)=\re(\Lambda_S'/\Lambda_S)(2+iT)=\sum_\rho a/(a^2+(T-\gamma)^2)$, $a=2-\beta\in[1,2]$ (F4). For $a\in[1,2]$, $a/(a^2+u^2)\ge1/(2+u^2)$, so $F(T)\ge\sum_\rho(2+(T-\gamma)^2)^{-1}$ with positive terms. For $\gamma>H$ let $d=\gamma-H'\ge H-H'$; by convexity of $T\mapsto T^{-2}$ and evenness of the kernel,
\[
\int_{H'}^\infty\frac{T^{-2}\,dT}{2+(T-\gamma)^2}\ge\int_{-d}^d\frac{(\gamma+u)^{-2}du}{2+u^2}\ge\gamma^{-2}\sqrt2\arctan(d/\sqrt2)\ge\frac{2}{c\gamma^2}.
\]
The zeros are symmetric under $\gamma\mapsto-\gamma$ and $F$ is even, so $\sum_{|\gamma|>H}2/\gamma^2\le2c\int_{H'}^\infty F(T)T^{-2}dT$. Now $F(T)$ is the sum of $\re(\gamma_S'/\gamma_S)(2+iT)$ and $\re(L_S'/L_S)(2+iT)$. (i) We have $\re(L_S'/L_S)(2+iT)=-\sum_n\Lambda_{\rm sym}(n)n^{-2}\cos(T\log n)$ and
\[
\Bigl|\int_{H'}^\infty\cos(aT)\,T^{-2}\,dT\Bigr|\le\frac2{aH'^2},
\]
so this part contributes at most $(2/H'^2)\sum_n|\Lambda_{\rm sym}(n)|/(n^2\log n)\le(6/H'^2)\sum_n\Lambda(n)/(n^2\log n)=6\log\zeta(2)/H'^2$. (ii) $(\gamma_S'/\gamma_S)(s)=-\frac32\log\pi+\frac12[\digamma(\frac{s+1}2)+\digamma(\frac{s+11}2)+\digamma(\frac{s+12}2)]$. For $z=x+iy$, $x,y>0$, Binet's formula for the digamma function $\digamma=\Gamma'/\Gamma$, $\digamma(z)=\log z-\frac1{2z}-2\int_0^\infty\frac{u\,du}{(u^2+z^2)(e^{2\pi u}-1)}$, $|u^2+z^2|\ge2xy$ and $\int_0^\infty u/(e^{2\pi u}-1)du=\frac1{24}$ give $\re\digamma(z)\le\log|z|+1/(24xy)$. With $z_j=x_j+iT/2$, $\log|z_j|\le\log(T/2)+2x_j^2/T^2$, so $\re(\gamma_S'/\gamma_S)(2+iT)\le\frac32\log(T/2\pi)+\sum_j[x_j^2/T^2+1/(24x_jT)]$; and $\int_{H'}^\infty\log(T/2\pi)T^{-2}dT=(\log(H'/2\pi)+1)/H'$. Finally $|\rho(\rho+2)|\ge\gamma^2$ and $|t^{1+\rho/2}|\le t^{3/2}$. The number $T(60,51.45)$ is evaluated in Arb (\file{tau-race/code/bounds.py}).
\end{proof}

(For orientation: $\sum_\rho|\rho|^{-2}=2(\Lambda_S'/\Lambda_S)(1)=0.198422$ and the $41$ zeros of Theorem~\ref{thm:4} give a true value of about $0.104$ for $\sum_{|\gamma|>60}2/\gamma^2$, if those zeros are on the line.)

\begin{lemma}[inputs on $\psi$ and $\theta$; unconditional]\label{lem:psi}
\textup{(a)} (Certified.) $\psi(y)\ge y-\sqrt2\sqrt y$ and $\theta(y)<y$ for $0<y\le Y_1:=10^7$. \textup{(b)} (F6.) For $y\ge10^7$: $(1-\varepsilon_1)y\le\theta(y)\le\psi(y)$ and $\theta(y)\le(1+\varepsilon_1)y$. Hence $\theta(y)\le c_\theta y$ for all $y>0$, and
\[
J_\zeta(t)\ge\tfrac23\bigl(1-\varepsilon_1\mathbf 1_{t>Y_1^2}\bigr)t^{3/2}-\tfrac45\sqrt2\,t^{5/4}\qquad(t>0).
\]
\end{lemma}

\begin{proof}
In (a), the supremum of $(y-\psi(y))/\sqrt y$ over $(0,10^7]$ is $\sqrt2$, approached as $y\to2^-$; between consecutive prime powers the quotient increases, so its supremum on each interval is its limit at the right end, which the program evaluates. For $\sqrt t\le Y_1$ integrate (a) in $J_\zeta(t)=\int_0^{\sqrt t}\psi(y)2y\,dy$; for $\sqrt t>Y_1$, $J_\zeta\ge\frac23Y_1^3-\frac45\sqrt2Y_1^{5/2}+\frac23(1-\varepsilon_1)(t^{3/2}-Y_1^3)\ge\frac23(1-\varepsilon_1)t^{3/2}-\frac45\sqrt2t^{5/4}$.
\end{proof}

\section{Proof of Theorem~\ref{thm:1}}\label{sec:proof1}

Let $t\ge1$ and use \eqref{eq:decomp}.
\begin{itemize}
\item $\psi_{1,\Delta}(t)\ge-\sigma_0^2t^{3/2}-a_0t$ by \eqref{eq:psi1bound} (this is where $\GRHD$ is used).
\item $J_\zeta(t)\ge\frac23(1-\varepsilon_1\mathbf 1_{t>10^{14}})t^{3/2}-\frac45\sqrt2t^{5/4}$ by Lemma~\ref{lem:psi}.
\item $J_S(t)\le Z_S(t)+2t^{1/2}+K_S$ by \eqref{eq:JSbound}. Under $\GRHS$, $|t^{1+\rho/2}|=t^{5/4}$ and $Z_S(t)\le C_St^{5/4}$ (Lemma~\ref{lem:lowzeros}). Under $\mathrm{RH}_S(60)$ only, the zeros with $|\gamma|\le60$ contribute at most $C_St^{5/4}$ and the others at most $T_{60}t^{3/2}$ (Lemma~\ref{lem:highzeros}).
\item Odd $k\ge3$: $|\int_0^t\Theta^{(k)}(u^{1/k})du|\le2c_\theta\frac k{k+1}t^{1+1/k}$, and the term vanishes unless $2^k\le t$. Keep $\frac34,\frac56$ for $k=3,5$; for odd $k\ge7$ use $\frac k{k+1}t^{1+1/k}\le t^{8/7}$; there are at most $\max(0,\frac12(\log t/\log2-5))$ such $k$. The total is $\le2c_\theta E_{\rm odd}(t)t^{3/2}$.
\end{itemize}
Adding up gives $\Theta_1(t)\ge L_A(t)t^{3/2}$ under $\GRHD$ and $\GRHS$, and $\Theta_1(t)\ge L_B(t)t^{3/2}$ under $\GRHD$ and $\mathrm{RH}_S(60)$, i.e.\ under $\GRHD$ by Theorem~\ref{thm:4}.

\emph{Monotonicity.} All of $t^{-1/2},t^{-1/4},t^{-1},t^{-3/2},t^{-1/6},t^{-3/10}$ decrease, and $t^{-5/14}(\log t/\log2-5)$ decreases for $\log t>5\log2+\frac{14}5$, i.e.\ $t>526.2$. So $L_A,L_B$ are nondecreasing on $[527,10^{14}]$ and on $(10^{14},\infty)$; since $L_B(10^6)\ge0.0274$ and $L_B(10^{14}+)\ge0.3612$ (Arb), $\Theta_1(t)>0$ for $t\ge10^6$.

\emph{Small $t$.} $\Theta_1$ is continuous and piecewise linear with nodes at the primes; $\Theta_1=0$ on $[0,2]$, $\Theta_1(t)=\lambda(2)\log2\,(t-2)<0$ on $(2,3]$ ($\lambda(2)=-24\cdot2^{-11/2}$), and on $[3,5]$, $\Theta_1(t)=\Theta(3)t-(2\lambda(2)\log2+3\lambda(3)\log3)$ with $\Theta(3)=0.2901\ldots>0$, which vanishes exactly at $t^*$ (certified enclosure $[4.266791956347541,4.266791956347542]$) and is positive on $(t^*,5]$. At every prime $5\le p\le9\,999\,991$ the inequality $\Theta_1(p)>0$ is certified (\file{tau-race/code/cert_small.py}: $664\,577$ primes, exact integers $\tau(p)$, $128$-bit balls). By linearity between nodes, $\Theta_1>0$ on $(t^*,9\,999\,991]$. \qed

(The ranges overlap from $10^6$ to $10^7$; $L_B$ is positive from $t\approx7.05\cdot10^5$, and $L_A$ from $t\approx1.112\cdot10^5$.)

\section{Proof of Theorem~\ref{thm:2}}\label{sec:proof2}

\begin{lemma}[second-order partial summation; standard, cf.\ Ingham]\label{lem:partial}
Let $x>2$, $2\le T\le x$ and $F\in C^2[2,x]$. Then $\sum_{p\le T}\lambda(p)\log p\,F(p)=F(T)\Theta(T)-F'(T)\Theta_1(T)+\int_2^T\Theta_1(t)F''(t)\,dt$.
\end{lemma}

\begin{proof}
Stieltjes integration by parts twice, with $\Theta_1(2)=0$.
\end{proof}

Take $F=h$, $h(t)=(x-t)/\log t$. Then $h(x)=0$, $h'(x)=-1/\log x$, and
\[
h''(t)=\frac{x+t}{t^2\log^2t}+\frac{2(x-t)}{t^2\log^3t}>0\qquad(1<t\le x).
\]
With $T=x$ this gives $xR(x)=\Theta_1(x)/\log x+\int_2^x\Theta_1h''$. For general $T\le x$, $\int_2^T\Theta_1h''=xA_1(T)+A_2(T)$ with
\[
A_1(T)=D_0(T)-\frac{\Theta(T)}{\log T}-\frac{\Theta_1(T)}{T\log^2T},\qquad A_2(T)=-D_1(T)+\frac{T\Theta(T)-\Theta_1(T)}{\log T}+\frac{\Theta_1(T)}{\log^2T}.
\]
Since $\Theta_1<0$ on $(2,t^*)$, the kernel term cannot simply be dropped; but its part up to $97$ is positive: certified, $A_1(97)\in[0.671972273685,0.671972273686]$ and $A_2(97)\in[10.326446649,10.326446650]$.

\begin{proof}[Proof of Theorem~\ref{thm:2}]
\emph{$x\ge97$.} By Theorem~\ref{thm:1}, $\Theta_1\ge0$ on $[97,x]$, and $h''>0$, so $xR(x)=\Theta_1(x)/\log x+xA_1(97)+A_2(97)+\int_{97}^x\Theta_1h''\ge\Theta_1(x)/\log x+0.6719x$. Hence $R(x)\ge\Theta_1(x)/(x\log x)+0.6719\ge L_B(x)\sqrt x/\log x+0.6719$, which is $>0$ (the middle term is $>0$ by Theorem~\ref{thm:1} even where $L_B<0$).

\emph{$x<97$.} On $[p,p')$ (consecutive primes), $R(x)=D_0(p)-D_1(p)/x$ is monotone, and $R$ is continuous. $R(p)>0$ is certified at every prime $7\le p\le9\,999\,991$, so $R>0$ on $[7,9\,999\,991]$. On $[5,7]$, $R(x)=D_0(5)-D_1(5)/x$ with $D_0(5)=0.7596\ldots$, $D_1(5)=4.1916\ldots$: increasing, with zero $x^*$. On $[3,5]$, $R$ is monotone with $R(3)=-0.1767\ldots$, $R(5)=-0.0787\ldots$; on $(2,3]$, $R(x)=\lambda(2)(1-2/x)<0$. Hence $R(x)>0\iff x>x^*$.

\emph{Constants.} For real $x\in[q,p]$ between consecutive primes $11\le q<p\le9\,999\,991$, $R(x)\log x/\sqrt x\ge\min(R(q),R(p))\log p/\sqrt p$ ($\log x/\sqrt x$ decreases for $x>e^2$). The program certifies that the minimum of these bounds is $\ge0.1718$ over all gaps with $q\ge11$, and $\ge0.7951$ over all gaps with $p>100$. For $x\ge9\cdot10^6$, $R(x)\log x/\sqrt x\ge L_B(9\cdot10^6)\ge0.1589$, resp.\ $\ge L_A(9\cdot10^6)\ge0.3435>\frac13$ under $\GRHS$.
\end{proof}

\begin{remark}[lead inheritance]
$D_0(x)=D_1(x)/x+R(x)$. So under $\GRHD$, for every $x>x^*$: $\sum_{p\le x}\tau(p)p^{-9/2}\ge0\Rightarrow\sum_{p\le x}\tau(p)p^{-11/2}>0$.
\end{remark}

\begin{remark}[what is not claimed]
The unsmoothed sum $D_0(x)=\sum_{p\le x}\lambda(p)$ is not one-signed ($D_0(17)<0$), and nothing pointwise is proved for the Koyama--Kurokawa weight $\tau(p)p^{-6}$. For other smoothings ($e^{-p/x}$, higher Riesz means) we expect that, under $\GRHD$, positivity for large $x$ follows from $\Theta_1>0$ on $(t^*,\infty)$ by Lemma~\ref{lem:partial}; this is not proved here, and the initial negative patch $(2,t^*)$ makes ``every $x$'' a separate finite computation for each weight, which we have not done.
\end{remark}

\begin{remark}[robustness of (F6)]\label{rem:robust}
Replace (F6) by the Rosser--Schoenfeld consequences quoted there, i.e.\ use $\varepsilon_1=1/\log10^7=0.062042\ldots$ and $c_\theta=1.01624$ in \eqref{eq:LA}. Then the explicit function of Theorem~\ref{thm:1} changes: $L_B(10^6)\ge0.0230$, $L_B(9\cdot10^6)\ge0.1562$, $L_B(10^{14}+)\ge0.3202$, $\lim L_B\ge0.3284$; $L_A(10^6)\ge0.2076$, $L_A(9\cdot10^6)\ge0.3408$, $\lim L_A\ge0.5130$. Everything else survives as stated: the sign statements of Theorems~\ref{thm:1} and \ref{thm:2}, Theorem~\ref{thm:3}, and the two constants of Theorem~\ref{thm:2} (as $0.1562>0.15$ and $0.3408>\frac13$; the finite range does not use (F6)).
\end{remark}

\section{Proof of Theorem~\ref{thm:4}}\label{sec:proof4}

\subsection{Evaluating \texorpdfstring{$\Lambda_S$}{Lambda\_S} rigorously}
$a_S(n)$ is multiplicative and rational, $a_S(p^k)=c\,a_S(p^{k-1})-c\,a_S(p^{k-2})+a_S(p^{k-3})$ with $c=\tau(p)^2/p^{11}-1$, and $|a_S(n)|\le d_3(n)\le n^2$ by (F1), where $d_3(n)$ is the number of ordered factorisations $n=abc$. Let
\[
K(z)=4z\int_0^\infty u^9\exp(-\pi z^2/u^2-2\pi u)\,du\quad(|\arg z|<\tfrac\pi4),\qquad\Phi(z)=\sum_{n\ge1}a_S(n)K(nz).
\]

\begin{lemma}\label{lem:V0}
\textup{(a)} $\int_0^\infty K(x)x^{s-1}dx=\gamma_S(s)$ for $\re s>-1$. \textup{(b)} For $|\arg z|=\vartheta<\frac\pi4$ and $0<\eta<1$: $|K(z)|\le C_\eta|z|\exp(-b_{\eta,\vartheta}|z|^{2/3})$ with $C_\eta=4\cdot9!/(2\pi\eta)^{10}$ and $b_{\eta,\vartheta}=3\pi(\cos2\vartheta\,(1-\eta)^2)^{1/3}$. \textup{(c)} $\Phi$ is analytic in the sector, $\Lambda_S(s)=\int_0^\infty\Phi(y)y^{s-1}dy$ for all $s\in\CC$, and $\Phi(1/z)=z\Phi(z)$.
\end{lemma}

\begin{proof}
(a) Integrate in $x$ first, then in $u$:
\[
\int_0^\infty x^se^{-\pi x^2/u^2}dx=\tfrac12\pi^{-(s+1)/2}u^{s+1}\Gamma\bigl(\tfrac{s+1}2\bigr),\qquad\int_0^\infty u^{s+10}e^{-2\pi u}du=(2\pi)^{-(s+11)}\Gamma(s+11).
\] (b) $|\exp(-\pi z^2/u^2)|=\exp(-\pi|z|^2\cos2\vartheta/u^2)$; by AM--GM, $\pi|z|^2\cos2\vartheta/u^2+2\pi(1-\eta)u\ge3\pi(\cos2\vartheta(1-\eta)^2|z|^2)^{1/3}$; and $\int_0^\infty u^9e^{-2\pi\eta u}du=9!/(2\pi\eta)^{10}$. (c) By (a), (b) and absolute convergence, $\int_0^\infty\Phi(y)y^{s-1}dy=\gamma_S(s)L_S(s)$ for $\re s>1$, and by Mellin inversion $\Phi(y)=\frac1{2\pi i}\int_{(2)}\Lambda_S(s)y^{-s}ds$. We move the line to $\re s=-1$. On $\re s=2$, $|L_S|\le\zeta(2)^3$; on $\re s=-1$, $|L_S(-1+iT)|=|\gamma_S(2-iT)/\gamma_S(-1+iT)|\,|L_S(2-iT)|\ll(|T|+2)^{9/2}$ by (F3) and Stirling. Since $L_S$ is entire of finite order (F3), Phragm\'en--Lindel\"of gives $L_S(\sigma+iT)\ll(|T|+2)^{9/2}$ uniformly for $-1\le\sigma\le2$, and Stirling gives $\gamma_S(\sigma+iT)\ll(|T|+2)^Ae^{-3\pi|T|/4}$ there for $|T|\ge1$. So $\Lambda_S(\sigma+iT)\to0$ exponentially, uniformly in the strip, the horizontal segments vanish, and $\Phi(y)=\frac1{2\pi i}\int_{(-1)}\Lambda_S(s)y^{-s}ds$ (on $\re s=-1$ the pole of $\gamma_S$ at $-1$ is cancelled by the zero of $L_S$). Substituting $s\mapsto1-s$ and using (F3) gives $\Phi(y)=y^{-1}\Phi(1/y)$ for $y>0$, hence in the sector by analytic continuation. Then $\Phi(y)$ decays faster than any power as $y\to0$ as well, and the Mellin integral converges for all $s$ and represents the entire function $\Lambda_S$.
\end{proof}

\begin{lemma}[trapezoidal rule and Poisson summation]\label{lem:V1}
For $h>0$ and $s\in\CC$,
\[
h\sum_{k\in\ZZ}\Phi(e^{kh})e^{skh}=\sum_{m\in\ZZ}\Lambda_S(s+2\pi im/h),
\]
and the left side equals $h[\Phi(1)+\sum_{k\ge1}\Phi(e^{kh})(e^{skh}+e^{(1-s)kh})]$.
\end{lemma}

\begin{proof}
$f(v)=\Phi(e^v)e^{sv}$ extends analytically to $|\im v|<\frac\pi4$ with decay faster than any exponential, uniformly in $|\im v|\le\frac\pi8$ (Lemma~\ref{lem:V0}); so Poisson summation applies, with $\hat f(\xi)=\Lambda_S(s-2\pi i\xi)$. The second form uses $\Phi(e^{-kh})=e^{kh}\Phi(e^{kh})$.
\end{proof}

\begin{lemma}[decay in the strip]\label{lem:V2}
For $0<\vartheta<\frac\pi4$, $-1\le\sigma\le2$ and $T\in\RR$: $|\Lambda_S(\sigma+iT)|\le M_\vartheta e^{-\vartheta|T|}$, $M_\vartheta:=2\int_1^\infty y\sup_\pm|\Phi(ye^{\pm i\vartheta})|dy$.
\end{lemma}

\begin{proof}
For $T>0$ rotate the ray in $\Lambda_S(s)=\int_0^\infty\Phi(z)z^{s-1}dz$ to $\arg z=\vartheta$, split at $|z|=1$ and apply $\Phi(z)=z^{-1}\Phi(1/z)$ on the inner part: $\Lambda_S(s)=e^{i\vartheta s}\int_1^\infty\Phi(ye^{i\vartheta})y^{s-1}dy+e^{i\vartheta(s-1)}\int_1^\infty\Phi(ye^{-i\vartheta})y^{-s}dy$. Here $|e^{i\vartheta s}|=|e^{i\vartheta(s-1)}|=e^{-\vartheta T}$, and $y^{\sigma-1},y^{-\sigma}\le y$ for $y\ge1$. For $T<0$ use $-\vartheta$.
\end{proof}

With $\vartheta=\pi/8$, $\eta=\frac15$ one has $b'=b_{\eta,\vartheta}=7.236>\frac92$, so $n\mapsto n^3e^{-b'(ny)^{2/3}}$ decreases for $n,y\ge1$, and summing gives $M_{\pi/8}\le109.74$.

\emph{Evaluation.} Take $h=2\pi/640$ and $k\le517$ ($e^{517h}>160$). For $-1\le\re s\le2$, $|\im s|\le73$,
\[
\Lambda_S(s)=h\Bigl[\Phi_0+\sum_{k=1}^{517}\Phi_k\bigl(e^{skh}+e^{(1-s)kh}\bigr)\Bigr]+E,\qquad|E|\le E_{\rm tail}+E_{\rm alias}\le3.2\cdot10^{-100}+4.4\cdot10^{-95},
\]
where the $\Phi_k$ are balls containing $\Phi(e^{kh})$, of radius $\le2.91\cdot10^{-84}$ (\file{tau-race/code/sym2_phi.py}). For $n\le\max(1,\lfloor160/y\rfloor)$, $K(ny)$ is computed by Arb's rigorous integrator on an exact interval $[u_{\rm lo},u_{\rm hi}]$, plus $\int_0^{u_{\rm lo}}\le u_{\rm lo}^{10}e^{-\pi x^2/u_{\rm lo}^2}$ and $\int_{u_{\rm hi}}^\infty\le\Gamma(10,2\pi u_{\rm hi})/(2\pi)^{10}$; the terms $n>N$ are bounded by $C_\eta y[(N+1)^3e^{-\beta(N+1)^{2/3}}+\frac32\Gamma(6,\beta(N+1)^{2/3})/\beta^6]$ ($\eta=\frac1{10}$, $\vartheta=0$, $\beta=b_\eta y^{2/3}$; first term plus integral, the summand being decreasing). $E_{\rm tail}=\sum_{k>517}h|\Phi(e^{kh})|\,2e^{2kh}$ is bounded by summing the bound $|\Phi(y)|\le C_\eta y(1+\frac32c_0)e^{-b_\eta y^{2/3}}$ ($y\ge1$, $c_0=\Gamma(6,b_\eta)e^{b_\eta}/b_\eta^6$; from Lemma~\ref{lem:V0}(b) and the monotonicity of $n^3e^{-b(ny)^{2/3}}$) beyond $y_0=e^{517h}$, and $E_{\rm alias}=\sum_{m\ne0}|\Lambda_S(s+640im)|\le2M_\vartheta e^{73\vartheta}e^{-640\vartheta}/(1-e^{-640\vartheta})$ by Lemmas~\ref{lem:V1}, \ref{lem:V2}. On the critical line $Z(t):=\Lambda_S(\frac12+it)=h[\Phi_0+\sum_k2\Phi_ke^{kh/2}\cos(tkh)]+E$ is real, and its computed enclosure has radius about $10^{-83}$ for $0\le t\le72$, against $|Z(60)|\approx9\cdot10^{-50}$ and $|Z(72)|\approx2\cdot10^{-61}$. Consistency: $L_S(2)$ from this formula is $0.80587520944866875961\ldots$, inside the Euler-product enclosure; $L_S(\frac12)=0.505549375222713515\ldots$ and $L_S(1)=0.631792945727883203\ldots$ agree with PARI to all digits printed.

\subsection{Zeros on the line}
\file{tau-race/code/sym2_verify.py} decides the sign of $Z$ on the grid $t=j/16$, $0\le j\le1152$: there are $55$ sign changes in $(0,72]$, $41$ of them in $(0,60]$. Each is bisected to an interval $[l_j,r_j]$ of length $2^{-38}$ with certified opposite signs at the ends, so there are zeros $\frac12+i\gamma_j$, $\gamma_j\in[l_j,r_j]$, $j=1,\dots,55$ (they agree with PARI's \texttt{lfunzeros} to all digits compared).

\subsection{Completeness}

\begin{lemma}\label{lem:V3}
Let $\gamma_1<\dots<\gamma_J$ be ordinates of zeros of $L_S$ on the critical line (one zero counted for each $j$), $w(u)=\frac32/(\frac94+u^2)$, and
\[
\mathrm{Rem}(T):=\re(\Lambda_S'/\Lambda_S)(2+iT)-\sum_{j}[w(T-\gamma_j)+w(T+\gamma_j)]\qquad(T\in\RR).
\]
Then \textup{(i)} $\mathrm{Rem}(T)\ge0$ and $|\mathrm{Rem}'(T)|\le\mathrm{Rem}(T)$; \textup{(ii)} if some zero $\rho_0=\beta_0+i\gamma_0$ of $L_S$ in the critical strip is not among the $\frac12\pm i\gamma_j$ (as a multiset), then $\mathrm{Rem}(\gamma_0)\ge\frac23$. Hence, if $\mathrm{Rem}(T)<\frac23$ for all $T\in[0,H]$, the zeros of $L_S$ with $0\le\im\rho\le H$ are exactly the $\frac12+i\gamma_j$ with $\gamma_j\le H$, each simple, and there is no zero with $\im\rho=0$.
\end{lemma}

\begin{proof}
By (F4), $\re(\Lambda_S'/\Lambda_S)(2+iT)=\sum_\rho f_\rho(T)$, $f_\rho(T)=a/(a^2+(T-\gamma)^2)$, $a=2-\beta\in[1,2]$, and $f_\rho=w(T\mp\gamma_j)$ for $\rho=\frac12\pm i\gamma_j$. So $\mathrm{Rem}$ is the sum of $f_\rho$ over the remaining zeros: (i) $\mathrm{Rem}\ge0$ and $|f_\rho'|=f_\rho\cdot2|u|/(a^2+u^2)\le f_\rho/a\le f_\rho$. (ii) If $\beta_0=\frac12$, $f_{\rho_0}(\gamma_0)=\frac23$. If $\beta_0\ne\frac12$, then $1-\beta_0+i\gamma_0$ is another remaining zero, and $\frac1{2-\beta_0}+\frac1{1+\beta_0}\ge\frac43$.
\end{proof}

\emph{Application.} The program computes, for $T=j/8$, $0\le j\le480$, an upper bound for $\mathrm{Rem}(T)$: the Euler product over $p\le2\cdot10^5$ with tail $\le3.03(\log P+1)/P$, $\gamma_S'/\gamma_S$ in Arb, and a lower bound for the found zeros' terms (all $55$, with $\gamma_j$ replaced by $[l_j,r_j]$). The result is $\max_j\mathrm{Rem}(j/8)\le0.185062$, attained at $T=60$; by (i), $\mathrm{Rem}(T)\le0.185062\,e^{1/16}\le0.196998<\frac23$ on $[0,60]$. Lemma~\ref{lem:V3} and the symmetry $\rho\mapsto\bar\rho$ give Theorem~\ref{thm:4}. \qed

($\mathrm{Rem}$ is small because the zeros up to $72$ are known; $\mathrm{Rem}(0)\approx0.0686$ is the contribution of all zeros above $72$. Mutation test: deleting any one of the zeros \#1, \#8, \#21, \#41 raises $\max\mathrm{Rem}$ to $0.743$, $0.741$, $0.758$, $0.842>\frac23$ (upper bounds computed by \file{tau-race/code/selftest.py} with $P=2\cdot10^4$ on a grid of step $\frac14$), so the test rejects.)

Lemma~\ref{lem:V3} is a variant of the method of Hiary--Ireland--Kyi \cite{HIK}: they use the positive sum $\sum_\rho(\rho-\delta)^{-1}$ at one real point $\delta\le0$, evaluated from $L'/L$ in the region of absolute convergence; here the positive partial-fraction sum is used at the points $2+iT$, $0\le T\le H$, with a Lipschitz bound.

\section{Proof of Theorem~\ref{thm:3}}\label{sec:proof3}

(a)$\Rightarrow$(b), (a)$\Rightarrow$(d) are Theorems~\ref{thm:2}, \ref{thm:1}; (b)$\Rightarrow$(c) and (d)$\Rightarrow$(e) are trivial.

\emph{(c)$\Rightarrow$(a)} (Landau). Let $A(x)=xR(x)=\sum_{p\le x}\lambda(p)(x-p)$; $|A(x)|\le2x^2$. For $\re s>1$, $\int_1^\infty A(x)x^{-s-2}dx=P(s)/(s(s+1))$ with $P(s)=\sum_p\lambda(p)p^{-s}$. From the Euler products, $\log L(s,\Delta)=\sum_p\sum_kc_k(p)/(kp^{ks})$ and $\log L_S(w)-\log\zeta(w)=\sum_p\sum_jc_{2j}(p)/(jp^{jw})$, so
\[
P(s)=\log L(s,\Delta)-\tfrac12\log L_S(2s)+\tfrac12\log\zeta(2s)+H(s),
\]
with $H(s)=\frac12\sum_p\sum_{j\ge2}c_{2j}(p)/(jp^{2js})-\sum_p\sum_{k\ge3}c_k(p)/(kp^{ks})$ absolutely convergent and analytic in $\re s>\frac13$. On the real segment $\sigma>\frac12$: $\zeta(2\sigma)>0$ is finite, $L_S(2\sigma)>0$ (Euler product), and $L(\sigma,\Delta)\ne0$, because for real $w$
\[
(2\pi)^{-w}\Gamma(w)L(w-\tfrac{11}2,\Delta)=\int_0^\infty\Delta(iy)y^{w-1}dy=\int_1^\infty\Delta(iy)(y^{w-1}+y^{11-w})dy>0,
\]
as $\Delta(iy)=e^{-2\pi y}\prod(1-e^{-2\pi ny})^{24}>0$. Hence $P(s)/(s(s+1))$ continues analytically to a neighbourhood of every real $\sigma>\frac12$. If $A(x)\ge0$ for $x\ge x_0$, Landau's theorem (the Mellin transform of an eventually one-signed function has a singularity at the real point of its abscissa of convergence; e.g.\ \cite[\S15.1]{MV}) shows that the integral converges and is analytic in $\re s>\frac12$. So $P$, and then $\log L(s,\Delta)=P+\frac12\log L_S(2s)-\frac12\log\zeta(2s)-H$, is analytic in $\re s>\frac12$, where $L_S(2s)\ne0$ and $\zeta(2s)\ne0,\infty$. So $L(s,\Delta)\ne0$ for $\re s>\frac12$, which is $\GRHD$ by the functional equation.

\emph{(e)$\Rightarrow$(a).} The same with $A$ replaced by $\Theta_1$: $\int_1^\infty\Theta_1(t)t^{-s-2}dt=P_1(s)/(s(s+1))$ with $P_1(s)=\sum_p\lambda(p)\log p\,p^{-s}=-(L'/L)(s,\Delta)+(L_S'/L_S)(2s)-(\zeta'/\zeta)(2s)+H_1(s)$ (Lemma~\ref{lem:decomp}), which is meromorphic in $\re s>\frac13$ with poles in $\re s>\frac12$ exactly at the zeros of $L(s,\Delta)$ there (simple, with residue $-\mu_\rho\ne0$), and no pole on the real segment $\sigma>\frac12$. Landau's theorem gives analyticity in $\re s>\frac12$, so there are no such zeros. \qed

\section{Computations}\label{sec:computations}

All programs are in \file{tau-race/code/}, outputs in \file{tau-race/data/}; Python 3.13 with python-flint 0.9.0 (Arb), PARI/GP for the independent checks. Exact $\tau(n)$ for $n\le10^7$ comes from $\Delta=q\prod(1-q^n)^{24}=q\,(\sum_{k\ge0}(-1)^k(2k+1)q^{k(k+1)/2})^8$ (Jacobi) with FLINT polynomial arithmetic (checked against $\tau\equiv\sigma_{11}\pmod{691}$, multiplicativity, the Hecke relation at $p^2$, PARI's \texttt{mfcoefs} and \texttt{ramanujantau}, and, by the second referee, against $(E_4^3-E_6^2)/1728$ at all $10^7$ values).

\emph{How $\sigma_0^2$ is certified.} With $w=s+\frac{11}2$, $\Lambda^*(w):=(2\pi)^{-w}\Gamma(w)L(w-\frac{11}2,\Delta)=\sum_n\tau(n)[E(w,n)+E(12-w,n)]$ with $E(w,n)=\int_1^\infty e^{-2\pi ny}y^{w-1}dy=(2\pi n)^{-w}\Gamma(w,2\pi n)$ (from $\Delta(i/y)=y^{12}\Delta(iy)$), and $\Lambda^{*\prime}(w)=\sum_n\tau(n)\int_1^\infty e^{-2\pi ny}\log y\,(y^{w-1}-y^{11-w})dy$. At $w=\frac{13}2$ we use $40$ terms, incomplete gamma functions in Arb, the derivative integrals by Arb's rigorous integrator on $[1,Y_n]$ plus an explicit tail, and a geometric bound for $n>40$ from $|\tau(n)|\le n^{6.5}$; then $\sigma_0^2=2\Lambda^{*\prime}/\Lambda^*(\frac{13}2)$. PARI's \texttt{lfun} (a different algorithm) gives $0.11225214047190852674558342625\ldots$, inside the enclosure. An independent audit summed over $3035$ zeros to height $2000$ plus a tail and matched $\sigma_0^2$ to $8\cdot10^{-11}$.

\emph{Mutation tests} (\file{tau-race/code/selftest.py}; each mutant must be rejected, and all are): sign-flipped $\lambda$, and $\tau(3),\tau(5),\tau(11)$ negated (the small-range checks fail); $\sigma_0^2:=0.7$ (no positivity threshold) and $H=12$ ($L_B$ never positive); one zero deleted ($\mathrm{Rem}\ge\frac23$); $\tau(2):=\tau(2)+1$ ($\sigma_0^2$ leaves its enclosure).


\part*{Part II. Excluded values}
\section{The equations}\label{sec:equations}

For a prime $p$ put $t=\tau(p)$. The Hecke relation $\tau(p^{j+1})=t\tau(p^j)-p^{11}\tau(p^{j-1})$ gives
\[
\tau(p^2)=t^2-p^{11},\qquad\tau(p^4)=t^4-3p^{11}t^2+p^{22}.
\]
So $\tau(p^4)=\alpha$ means that $(U,V)=(p^{11},t^2)$ is an integer solution of
\begin{equation}\label{eq:quad}
U^2-3UV+V^2=\alpha
\end{equation}
in which $U$ is an eleventh power and $V$ is a square, and $\tau(p^2)=c$ means $t^2-p^{11}=c$. Theorem~\ref{thm:p4} is the statement that \eqref{eq:quad} with $\alpha\in\{461,-599\}$ has no integer solution with $U=x^{11}$ and $V=y^2$.

\section{\texorpdfstring{$\tau(p^4)\ne461,-599$}{tau(p\^{}4) is not 461, -599}}\label{sec:p4}

\subsection{All solutions lie on finitely many recurrence sequences}
Fix $\alpha\ne0$ and let $S$ be the set of integer solutions of \eqref{eq:quad}. The map $\mathsf T(U,V)=(V,3V-U)$ is a bijection of $S$, with inverse $(U,V)\mapsto(3U-V,U)$, because $V^2-3V(3V-U)+(3V-U)^2=U^2-3UV+V^2$. For $(U,V)\in S$ put $Z_0=U$, $Z_1=V$ and $Z_{n+1}=3Z_n-Z_{n-1}$ for all $n\in\ZZ$; then $\mathsf T^n(U,V)=(Z_n,Z_{n+1})$.

\begin{lemma}[seed bound]\label{lem:seed}
Every $\mathsf T$-orbit in $S$ contains a pair $(U,V)$ with $U^2\le\alpha/2$ if $\alpha>0$, and with $U^2\le|\alpha|$ if $\alpha<0$.
\end{lemma}

\begin{proof}
Let $(Z_n)$ be the sequence of an orbit and choose $m$ with $|Z_m|$ minimal (the values are integers). Since $Z_{m+1}+Z_{m-1}=3Z_m$ and $(Z_m,Z_{m+1})\in S$, we get $Z_{m+1}Z_{m-1}=Z_{m+1}(3Z_m-Z_{m+1})=Z_m^2-\alpha$, and by minimality $|Z_{m+1}Z_{m-1}|\ge Z_m^2$. If $\alpha>0$: $Z_m^2\ge\alpha$ would give $0\le Z_m^2-\alpha<Z_m^2$, a contradiction; so $Z_m^2<\alpha$, and then $\alpha-Z_m^2\ge Z_m^2$ (this includes $Z_m=0$). If $\alpha<0$: $Z_{m+1}Z_{m-1}=Z_m^2+|\alpha|>0$, so $Z_m\ne0$ (otherwise the two neighbours would have sum $0$ and positive product), and both neighbours have the sign of $Z_m$. After changing the sign of the whole sequence, $a=Z_{m-1}$ and $b=Z_{m+1}$ satisfy $a,b\ge z:=Z_m>0$ and $a+b=3z$, hence $a,b\in[z,2z]$ and $ab=a(3z-a)\ge2z^2$. So $z^2+|\alpha|\ge2z^2$.
\end{proof}

So $S$ is the union of the orbits of the finitely many seeds $(U,V)\in S$ with $|U|\le B_0$, $B_0=\lfloor\sqrt{\alpha/2}\rfloor$ resp.\ $\lfloor\sqrt{|\alpha|}\rfloor$. For such $U$, $V=(3U\pm s)/2$ with $s^2=5U^2+4\alpha$.
\begin{itemize}
\item $\alpha=461$: $B_0=15$, and $5U^2+1844$ is a square only for $U=\pm1$ ($43^2$). Seeds: $(1,23)$, $(1,-20)$, $(-1,20)$, $(-1,-23)$.
\item $\alpha=-599$: $B_0=24$, and $5U^2-2396$ is a square only for $U=\pm24$ ($22^2$). Seeds: $(24,47)$, $(24,25)$, $(-24,-25)$, $(-24,-47)$.
\end{itemize}
Hence, if $x^{22}-3x^{11}y^2+y^4=\alpha$, then for one of the four seeds and some $n\in\ZZ$,
\begin{equation}\label{eq:target}
Z_n=x^{11}\quad\text{and}\quad Z_{n+1}=y^2.
\end{equation}

\subsection{Periodicity and the sieve}
Let $A=\begin{pmatrix}0&1\\-1&3\end{pmatrix}$, so that $(Z_n,Z_{n+1})^{\mathsf t}=A^n(Z_0,Z_1)^{\mathsf t}$ for all $n\in\ZZ$ ($\det A=1$). If $A^M\equiv I\pmod q$ for a prime $q$, then $n\mapsto(Z_n,Z_{n+1})\bmod q$ has period dividing $M$, for every seed and in both directions. The orders $T_q$ of $A$ modulo $q$ used are
\[
\begin{array}{c|cccccc}
q&3&7&43&89&199&967\\\hline T_q&4&8&44&22&11&88
\end{array}
\]
all dividing $88$. If \eqref{eq:target} holds then, for every prime $q$, $Z_n\bmod q\in P_q:=\{x^{11}\bmod q\}$ and $Z_{n+1}\bmod q\in Q_q:=\{y^2\bmod q\}$ ($0$ included in both). A class $n\bmod M$ is \emph{killed} by $q$ (with $T_q\mid M$) if $Z_n\bmod q\notin P_q$ or $Z_{n+1}\bmod q\notin Q_q$; this depends only on $n\bmod M$.

\begin{proof}[Proof of Theorem~\ref{thm:p4}]
For $\alpha=-599$, take $M=22$ and $q=89$ alone ($89=8\cdot11+1$, so $P_{89}$ has $9$ elements): for each of the four seeds all $22$ classes are killed. For $\alpha=461$, take $M=88$ and $q\in\{3,7,43,89,199,967\}$: for each of the four seeds all $88$ classes are killed. So \eqref{eq:target} is impossible. The certificate \file{tau-p4/certificate_tau_p4.txt} lists, for every seed and every class, a witness $(q,\text{code},Z_n\bmod q,Z_{n+1}\bmod q)$.
\end{proof}

\begin{remark}[which condition does the work]\label{rem:which}
For $\alpha=-599$ the eleventh-power condition alone suffices: with $M=22$ and $q\in\{89,199\}$ no class has $Z_n\in P_{89}\cap P_{199}$ simultaneously. So $u^2=5x^{22}-2396$ has no integer point, and this case falls to the method of \cite[Lemma~2.15]{Goto} (due to Dembner and Jain \cite{DJ}); we do not know why arXiv:2405.16723v1 lists $-599$ as open. For $\alpha=461$ the square condition is essential: the sequence contains $Z_0=\pm1$, an eleventh power, so no congruence on $Z_n$ alone can work, and with the eleventh-power condition only the classes $n\equiv0,11\pmod{22}$ survive on every seed. The neighbours of $\pm1$ are $23,-20$ (resp.\ $-23,20$), none a square, and the square condition kills those classes. The six witness primes for $461$ are not minimal ($43$ or $89$ can be dropped).
\end{remark}

\emph{Verification.} Three implementations agree: a Python/numpy sieve with modulus $55440$ and all primes up to $10^5$ whose period divides $55440$ (zero survivors; planted solutions $(U,V)=(2^{11},9)$ and $(3^{11},25)$ survive in exactly one class, as they must); a separately written PARI/GP generator using Euler's criterion and Kronecker symbols, which regenerates seeds, periods and the full witness table; and an independent Python checker with exact integers and brute-force residue sets and seeds, which re-verifies all $1376$ class lines of the certificate. The referee confirmed that every solution of \eqref{eq:quad} with $|U|\le2\cdot10^6$ lies on the listed orbits.

\section{\texorpdfstring{$\tau(p^2)\ne601$}{tau(p\^{}2) is not 601}, and the status of 463}\label{sec:p2}

Let $c\in\{463,601\}$ and suppose $t^2-p^{11}=c$ with $t=\tau(p)$. Then:
\begin{itemize}
\item $p$ is odd, since $\tau(4)=-1472$; and $p\ne c$, since $c^{11}+c$ is divisible by $c$ exactly once and so is not a square;
\item $p\ne23$ ($23^{11}+c$ is not a square) and $t\equiv2\pmod{23}$. Indeed $c\equiv3\pmod{23}$ for both values, and by Wilton's congruence \cite{SwD} $\tau(p)\equiv0\pmod{23}$ if $(\frac p{23})=-1$, and $\tau(p)\in\{2,-1\}\pmod{23}$ if $(\frac p{23})=1$, while $p^{11}\equiv(\frac p{23})$. In the first case $t^2-p^{11}\equiv1$; in the second $t^2\equiv4$, so $t\equiv\pm2$, and $-2\notin\{2,-1\}$.
\end{itemize}
Let $K=\QQ(\sqrt c)$. Its class number is $1$ in both cases (PARI \texttt{bnfinit} with \texttt{bnfcertify}; for $c=601$ also by the analytic class number formula, computed by the referee). Fundamental units:
\begin{align*}
c=463:&\quad\varepsilon_0=247512720456368+11502891625161\sqrt{463}\quad(\text{norm }+1),\\
c=601:&\quad\varepsilon_0=139468303679532+5689030769845\sqrt{601}\quad(\text{norm }-1),
\end{align*}
recomputed from continued fractions. A prime ideal dividing both $t\pm\sqrt c$ divides $2\sqrt c$ and $p$, which is impossible as $p$ is odd and $p\ne c$. So $(t+\sqrt c)$ is the eleventh power of an ideal, which is principal; absorbing $-1=(-1)^{11}$ and $\varepsilon_0^{11}$,
\begin{equation}\label{eq:descent}
t+\sqrt c=\varepsilon_0^kg^{11},\qquad g\in\cO_K,\ k\in\{0,1,\dots,10\},\ N(g)=\pm p .
\end{equation}

\emph{Local test.} Let $q\ne p$ be a prime. Then $g$ is a unit modulo $q$, so $(t+\sqrt c)\varepsilon_0^{-k}$ is an eleventh power in $(\cO_K/q)^\times$. Let $A(q,k)$ be the set of residues $t\bmod q$ for which this holds (at $q=23$ only $t\equiv2$ is admitted). If $A(q,k)=\emptyset$ then \eqref{eq:descent} is impossible for that $k$. For every $q$ used, $q^{11}+c$ is not a square, so $q\ne p$. For odd $q$, $\cO_K/q\cong\FF_q[y]/(y^2-c)$, so half-integral $g$ are covered.

\begin{proof}[Proof of Theorem~\ref{thm:p2} for $601$]
Both $23$ and $67$ split in $\QQ(\sqrt{601})$ ($601\equiv3\equiv7^2\pmod{23}$, $601\equiv-2\pmod{67}$, a square), the eleventh powers in $\FF_{23}^\times$ are $\pm1$, and $67\equiv1\pmod{11}$. With $\varepsilon_0$ as orientation, $A(23,k)=\emptyset$ for every $k$ except $k=7$, and $A(67,k)=\emptyset$ exactly for $k\in\{0,1,4,7,10\}$. So every $k\in\ZZ/11$ is excluded, and there is no prime $p$ with $\tau(p^2)=601$.
\end{proof}

This uses only the Hecke relation, Wilton's congruence, $11\nmid h(\QQ(\sqrt{601}))$, the fundamental unit, and finite computations modulo $23$ and $67$: no Thue equation and no GRH. Wilton's congruence is load-bearing: without the restriction $t\equiv2$, the classes $k=3,8$ survive every prime $q<700$. (The class $k=7$ is also killed by $q=89$.)

\subsection{The case 463}\label{sec:463}
For $c=463$, $A(23,k)=\emptyset$ for every $k$ except $k=9$, and the class $k=9$ survives every local test tried (all $q<700$ with $11\mid q^2-1$, all split $q\equiv1\pmod{11}$ up to $20000$, and the $11$-adic test), so a global argument is needed. Here $\cO_K=\ZZ[\sqrt{463}]$. Write $g=aw_1+bw_2$ with the $\ZZ$-basis $w_1=624+29\sqrt{463}$, $w_2=1829+85\sqrt{463}$ (determinant $-1$, chosen by LLL to make the coefficients small). Comparing coefficients of $\sqrt{463}$ in $t+\sqrt{463}=\varepsilon_0^{-2}g^{11}$ gives the Thue equation
\begin{multline*}
Q(a,b):=2069a^{11}+143a^{10}b+1072720a^9b^2-899250a^8b^3+61199160a^7b^4\\-48219864a^6b^5+821606016a^5b^6-492420720a^4b^7+2836978320a^3b^8\\-1043992400a^2b^9+1842900224ab^{10}-226837664b^{11}=1 .
\end{multline*}
$Q(x,1)$ is irreducible with one real root; its field has signature $(1,5)$ and discriminant $-2^{10}11^{11}463^5$. (The referee verified the form exactly; since $N(\varepsilon_0)=+1$ the classes $k$ and $-k$ are conjugate, so the orientation does not matter.)

\emph{The computation.} PARI/GP 2.15.4 \cite{PARI}, \texttt{thue} with an uncertified \texttt{thueinit} (flag $0$), finds no solution of $Q(a,b)=\pm1$, for three different inputs ($Q(x,1)$, $Q(1,x)$, and $Q(a+b,b)$), a few seconds each; a brute-force search finds none with $|a|,|b|\le500$; the tentative class number of the degree-$11$ field is $1$, with tentative regulator $7.6\cdot10^9$. PARI's documentation states that an uncertified result ``is conditional on the GRH if $a\ne\pm1$ and~[\dots]'', and that results computed with such a \texttt{tnf} are likewise conditional ``except for a right-hand side of $\pm1$''; here the right-hand side is $\pm1$. We have not proved this exemption. For right-hand side $\pm1$ the class group is irrelevant, but the answer depends on having a full system of fundamental units, which an uncertified \texttt{bnfinit} does not prove. A certified run (\texttt{thueinit}$(P,1)$, which calls \texttt{bnfcertify} on the degree-$11$ field) did not finish in two attempts totalling about $2$\,h\,$47$\,min (\file{tau-p2/thue463_certified_run.txt} records the first; the second is recorded in \file{tau-p2/TAU2.md}, \S5); the Minkowski bound is about $3.7\cdot10^{10}$, and a certification of the units alone may be cheaper. \emph{Thus the exclusion of $\tau(p^2)=463$ rests on this uncertified PARI computation.} A positive control ($c=94$, where $421^2-94=3^{11}$) recovers the known solution in the expected unit class; it tests the set-up, not the completeness of \texttt{thue}. A direct search found no odd $x<10^7$ with $x^{11}+463$ or $x^{11}+601$ a square.

\section{Consequences for odd primes below 1000}\label{sec:goto}

\subsection{The reduction}
Let $\ell$ be an odd prime with $\ell\nmid\tau(\ell)$ (for $\ell<1000$ this excludes only $3,5,7$) and suppose $\tau(n)=\pm\ell$. By \cite[Theorem~1.1]{BCOT} (proved in its Section~3, using Proposition~3.1 there and $|\tau(m)|\ne1$ for $m>1$), $n=p^{d-1}$ with $p$ and $d$ odd primes and $d\mid\ell(\ell^2-1)$. If $p\mid\tau(p)$ then $p\mid\tau(p^j)$ for all $j\ge1$ by the Hecke recursion, so $p=\ell$, contradicting $\ell\nmid\tau(\ell)$. For primes $p\nmid\tau(p)$ the sequence $u_m=\tau(p^{m-1})$ is a Lucas sequence, and \cite[Section~5]{BCOT} shows (from the classification of defective Lucas numbers, \cite[Lemma~2.1]{BCOT}, \cite{BHV}) that it has no defective term equal to $\pm1$ or $\pm\ell$; so $\ell$ is a primitive prime divisor of $u_d$, $\ell\nmid(\alpha-\beta)^2p^{11}$, and $d\mid\ell-(\frac{(\alpha-\beta)^2}\ell)$, in particular $d\mid\ell^2-1$ and $d\ne\ell$ (this step is written out in \cite[\S3.1]{Goto}). We rely on these published statements and did not re-derive the analysis of defective Lucas numbers.

\subsection{Congruences}
For $n=p^{d-1}$ with $p$ an odd prime not in $\{3,5,7,23,691\}$ (so $n\equiv1\pmod8$ and $n\equiv1\pmod3$) we use the classical congruences \cite{SwD}: $\tau(n)\equiv\sigma_{11}(n)$ modulo $691$ and modulo $2^{11}$; $\tau(n)\equiv n^{-610}\sigma_{1231}(n)\pmod{3^6}$; $\tau(n)\equiv n^{-30}\sigma_{71}(n)\pmod{5^3}$; $\tau(n)\equiv n\sigma_3(n)\pmod7$; and Wilton's congruence modulo $23$ followed by the Hecke recurrence. Each was re-checked numerically for all $n\le3000$ in the stated class (and by the referee for $n\le4000$), and the resulting residue sets against exact values $\tau(p^{d-1})$ for $p\le113$, $d\in\{3,5,7,11,13,23\}$; for $p\in\{3,5,7,23,691\}$ the values $\tau(p^{d-1})\ne\pm\ell$ were checked directly. Among all pairs $(\pm\ell,d)$ with $11\le\ell<1000$ prime and $d$ an odd prime dividing $\ell^2-1$, the congruences exclude all but seventeen:
\[
\begin{array}{ll}
d=3:&-137,\ -227,\ 463,\ 601,\ 691,\ -827\\
d=5:&-139,\ -179,\ 281,\ 461,\ -599,\ -919\\
d=23:&-139,\ -367;\qquad d=59:\ 827;\qquad d=83:\ -829;\qquad d=107:\ -643.
\end{array}
\]
(The same list was obtained by two further independent computations.) Table~\ref{tab:pairs} shows which moduli exclude each pair for $\pm461$ and $\pm599$.

\begin{table}[ht]\small
\caption{Pairs $(\varepsilon\ell,d)$, $d\mid\ell^2-1$, for $\ell=461,599$ ($461^2-1=2^3\cdot3\cdot5\cdot7\cdot11\cdot23$, $599^2-1=2^4\cdot3\cdot5^2\cdot13\cdot23$).}\label{tab:pairs}
\begin{tabular}{rr|l}\toprule
value&$d$&excluded modulo\\\midrule
$461$&$3$&$3^6$\\
$461$&$5$&\textbf{none} (Theorem~\ref{thm:p4})\\
$461$&$7$&$691,\ 3^6,\ 7$\\
$461$&$11,23$&$7$\\
$-461$&$3$&$691,\ 5^3,\ 23$\\
$-461$&$5$&$691,\ 2^{11},\ 5^3$\\
$-461$&$7$&$3^6,\ 5^3,\ 23$\\
$-461$&$11,23$&$5^3$\\
$599$&$3$&$3^6,\ 5^3$\\
$599$&$5$&$2^{11},\ 5^3$\\
$599$&$13$&$691,\ 2^{11},\ 3^6,\ 5^3,\ 7$\\
$599$&$23$&$5^3$\\
$-599$&$3$&$691,\ 23$\\
$-599$&$5$&\textbf{none} (Theorem~\ref{thm:p4})\\
$-599$&$13$&$3^6,\ 7,\ 23$\\
$-599$&$23$&$7$\\\bottomrule
\end{tabular}
\end{table}

\begin{proof}[Proof of Corollary~\ref{cor:values}]
For $\ell\in\{461,599\}$: both are ordinary ($\tau(461)\equiv101\pmod{461}$, $\tau(599)\equiv280\pmod{599}$), so the reduction applies, and Table~\ref{tab:pairs} with Theorem~\ref{thm:p4} excludes every pair. For $601$ ($601^2-1=2^4\cdot3\cdot5^2\cdot7\cdot43$, and $601$ is ordinary): the only surviving pair is $(601,3)$, i.e.\ $\tau(p^2)=601$, excluded by Theorem~\ref{thm:p2}; $-601$ does not occur in the list. For $\pm463$: $-463$ does not occur in the list, and $(463,3)$ is the case of Section~\ref{sec:463}.
\end{proof}

\subsection{The remaining pairs, and Goto's arXiv version}
\begin{proposition}\label{prop:minusD}
$\tau(p^2)\ne-137,-227,-827$ for every prime $p$.
\end{proposition}

\begin{proof}
For $D=827$: $t^2+827=p^{11}$ with $p$ odd and $p\ne827$ ($827^{11}-827$ is divisible by $827$ exactly once). The class number of $\QQ(\sqrt{-827})$ is $7$, prime to $11$, and the units are $\pm1$, so $t+\sqrt{-827}=((a+b\sqrt{-827})/2)^{11}$; comparing coefficients of $\sqrt{-827}$ gives $b\,F(a,b)=2^{11}$ for an explicit binary form $F$, so $b=\pm2^j$, and a finite check shows there is no solution. $D=227$ is identical (class number $5$). For $D=137$ the ring is $\ZZ[\sqrt{-137}]$ (discriminant $-548$, class number $8$), $t+\sqrt{-137}=\pm(a+b\sqrt{-137})^{11}$, and $b\,F(a,b)=\pm1$ has no solution. The class numbers were certified with \texttt{bnfcertify} and confirmed by the referee by enumerating reduced forms.
\end{proof}

The remaining pairs: $(281,5)$ and $(-919,5)$ are excluded by \cite[Lemma~2.15]{Goto} (and again by our sieve with $M=22$); $(-139,5)$, $(-179,5)$ likewise by the sieve; $(691,3)$ by \cite[Theorem~1.2(1)]{BCOT}; $(-139,23)$, $(-367,23)$ by the Thue equations $F_{22}(X,Y)=-139,-367$ over $\QQ(\zeta_{23})^+$, solved by PARI with a certified \texttt{thueinit}; and $(827,59)$, $(-829,83)$, $(-643,107)$ by the Thue equations of degrees $29$, $41$, $53$ of \cite[Lemma~2.14]{Goto}, which we re-ran with an uncertified \texttt{thueinit} (tentative class number $1$ in all three fields; by PARI's documented criterion such results are unconditional, a criterion we have not proved). For $\ell=3,5,7$ the result $\tau(n)\notin\{\pm1,\pm3,\pm5,\pm7\}$ is \cite[Theorem~1.1]{BCO}. In sum:

\begin{corollary}\label{cor:all}
For every odd prime $\ell<1000$ and every $n\ge1$, $\tau(n)\ne\pm\ell$, subject to: the cited results of \cite{BCO,BCOT}; the classical congruences; the uncertified (possibly GRH-dependent, since PARI's criterion refers to a tentative class number computed assuming GRH) PARI computations for $+827$, $-829$, $-643$ described above; and, for $\tau(n)=463$, the PARI computation of Section~\ref{sec:463}.
\end{corollary}

\emph{Goto's arXiv version.} In arXiv:2405.16723v1 the proof of Theorem~0.3 lists the pairs not excluded by Lemma~A.1 there, and the list contains no pair with $d=3$. The pair $(601,3)$ passes every congruence of Lemma~A.1 as printed ($601\equiv1\pmod{12}$, $\equiv1\pmod5$, $\equiv6\pmod7$, $\equiv3\pmod{23}$), so its omission is an error in that list. Separately, the mod-$12$ line ``$\tau(p^{d-1})\equiv0,1,d\pmod{12}$'' is false for $d=3$: $\tau(5^2)=-25499225\equiv7\pmod{12}$, the true set for $d=3$ being $\{1,3,7,9\}$ (one of our referees found it false for $d=11$ too); it wrongly excludes $463\equiv7\pmod{12}$. The pair $(-827,3)$ is not excluded by any of our congruences either. So v1 does not exclude $\tau(p^2)=463,601$, nor $-827$. Goto's paper is the source of the problem and of the methods for $281$, $-919$ and the large-$d$ Thue equations; the published version \cite{Goto} may differ from v1, and we have not seen it.

\begin{proposition}[two cases of Hanada--Madhukara]\label{prop:HM}
Let $\tau_{2k}$ denote the coefficients of the normalised level-$1$ eigenform of weight $2k\in\{16,26\}$, and $e=2k-1$. Then $\tau_{2k}(p^4)\notin\{33,55\}$ for every prime $p$.
\end{proposition}

\begin{proof}
As in Section~\ref{sec:equations}, $\tau_{2k}(p^4)=U^2-3UV+V^2$ with $(U,V)=(p^e,\tau_{2k}(p)^2)$. For $33$: $U^2-3UV+V^2=33$ has no integer solution, since modulo $3$ it forces $3\mid U$, $3\mid V$, and then $9\mid33$. For $55$: with $e=15$ ($M=40$) and $e=25$ ($M=84$) the sieve of Section~\ref{sec:p4} kills every class (for instance with $q\in\{3,5,7,11,41\}$, resp.\ $q\in\{3,13,29,83,211,281,421,1427\}$, as found by a referee) except $n\equiv0\pmod M$ for the seed $(Z_0,Z_1)=(1,9)$, which is the genuine solution $1-27+81=55$ with $U=1$. Suppose $Z_n=x^e$ and $Z_{n+1}=y^2$ with $n\ne0$, $n\equiv0\pmod4$, and write $n=2Kk$ with $K=2^a$, $a\ge1$, $k$ odd. With $V_K=\lambda^K+\lambda^{-K}$, $\lambda=(3+\sqrt5)/2$ ($V_1=3$, $V_{2K}=V_K^2-2$), one has $Z_{m+2K}=V_KZ_{m+K}-Z_m$, hence $Z_{1+2Kk}\equiv-Z_1=-9\pmod{V_K}$. From $V_2=7$ and the recursion, $V_K\equiv7\pmod8$ and $V_K\equiv1$ or $2\pmod3$ for all $K=2^a$, $a\ge1$; so the Jacobi symbol $(\frac{-9}{V_K})=(\frac{-1}{V_K})=-1$, while $y^2\equiv-9\pmod{V_K}$ with $\gcd(y,V_K)=1$ would make it $+1$ (the argument of Cohn \cite{Cohn}). So the only solutions have $x^e=1$, and $x$ is not a prime.
\end{proof}

In \cite[Remark~6.2]{HM} the integer points of these four curves are left undetermined; by the proof in \cite{HM}, the values $33$ and $55$ were left out of its lists for weights $16$ and $26$ only because of the hyperelliptic curves $Y^2=5X^{2e}+4c$, $c\in\{33,55\}$; so they can now be included there, \emph{provided} the other steps of \cite{HM} for these $c$ are as stated (we did not re-verify them). Proposition~\ref{prop:HM} determines only the solutions with $V$ a square; it does not list all integer points of those curves.


\section{Review, status and reproducibility}\label{sec:review}

\emph{Part I.} The working note \file{tau-race/TAU-RACE.md} was read by three independent AI agents. (1) A referee instructed to find the error re-derived Lemmas~\ref{lem:decomp}--\ref{lem:partial}, the assembly of Section~\ref{sec:proof1}, Lemmas~\ref{lem:V0}--\ref{lem:V3} and Section~\ref{sec:proof3}, reproduced all certificates byte for byte, and re-checked the identity \eqref{eq:decomp}, the explicit formula for $J_S$, $T(60,51.45)$, the completeness test and the zero counts with independent PARI computations; verdict ``repairable gaps; no fatal error'' (\file{tau-race/REFEREE-REPORT.md}); all repairs were applied. (2) An independent numerical audit, in floating point, with its own PARI and mpmath code (\file{tau-race/AUDIT-NUMERICS.md}), including a floating-point argument-principle zero count giving $41$ zeros up to height $60$ and $55$ up to $72$: six items, all pass. (3) A second referee with its own code (\file{tau-race/REFEREE-2-TAU.md}), which recomputed $\tau(n)$ for $n\le10^7$ from $(E_4^3-E_6^2)/1728$: no fatal error and no wrong step in Theorems~\ref{thm:1}--\ref{thm:4}, each ``complete modulo listed gaps''. Its three gaps were citation and presentation gaps and are repaired here: the functional equation of $\Lambda_S$ is quoted from sources actually read (Zagier \cite{Zagier}, Balkanova--Frolenkov \cite{BF}; (F3)); the contour shift of Lemma~\ref{lem:V0}(c) is proved; and (F6) is quoted from the texts of \cite{BKLNW,Dusart}, with Remark~\ref{rem:robust} stating exactly which constants change under the Rosser--Schoenfeld fallback. These repairs were made after that report and were not refereed again. Not consulted by anyone: the originals of Shimura, Gelbart--Jacquet, Rankin, Rosser--Schoenfeld, Deligne, Hecke and Swinnerton-Dyer, the finite order of $L_S$ (which is cited), and the books \cite{IK,MV}, whose numbered statements are quoted from memory. All reads concern the working notes, not this manuscript. No second rigorous implementation of the $\Phi$-table or of the completeness test was written; Theorem~\ref{thm:4} rests on Arb, the two programs and (F3), (F4).

\emph{Part II.} \file{tau-p4/PROOF.md} was read by two independent agents with their own scripts (the first read is summarised in \file{tau-p4/PROOF.md}, \S7; the second report is \file{tau-p4/REFEREE-2-TAU4.md}): Theorem~\ref{thm:p4} ``proof complete''; the deduction of Corollary~\ref{cor:values} for $\pm461,\pm599$ ``complete, modulo the cited literature''. The first read confirmed the seed bound (also by a stress test over all $0<|\alpha|\le700$), the seeds, the periods, zero survivors with its own sieve, the congruences, the seventeen surviving pairs, the class-number arguments of Proposition~\ref{prop:minusD}, Proposition~\ref{prop:HM}, and the quotations from \cite{Goto,BCO,BCOT,BGPS,HM}. \file{tau-p2/TAU2.md} was read by an independent referee (\file{tau-p2/REFEREE-2-TAU2.md}): the case $601$ ``proof complete''; the labelling of the case $463$ honest. Applied corrections: the reduction is cited to \cite{BCOT} directly (Theorem~1.1, Proposition~3.1, Section~5); the per-pair table (Table~\ref{tab:pairs}) was added; the cases resting on uncertified PARI computations are labelled as such.

\emph{Programs.} Part I: \file{tau-race/code/} (\file{tau_table.py}, \file{cert_constants.py}, \file{cert_small.py}, \file{bounds.py}, \file{sym2_phi.py}, \file{sym2_verify.py}, \file{selftest.py}, \file{independent_check.gp}), outputs in \file{tau-race/data/} with SHA-256 hashes in \file{tau-race/TAU-RACE.md}, \S7. Total running time a few minutes on four cores. Part II: \file{tau-p4/tau_p4.gp} (generator, $12$\,s), \file{tau-p4/certificate_tau_p4.txt}, \file{tau-p4/check_certificate.py} (independent checker), \file{tau-p4/sieve_tau_p4.py}, \file{tau-p4/thue_check.gp}; \file{tau-p2/local_k.gp}, \file{tau-p2/local_k_check.py}, \file{tau-p2/thue463.gp}, \file{tau-p2/thue463_runs.gp}, \file{tau-p2/control94.gp}.

\section*{Acknowledgements}
This paper was produced by AI (Anthropic's Claude models), directed by the author: the statements, proofs, programs, computations and literature searches were generated by AI. It was checked by independent AI referees, as described in Section~\ref{sec:review}, and has not yet been refereed by a human expert. The author thanks the developers of FLINT/Arb, python-flint and PARI/GP.

\begin{thebibliography}{99}\small

\bibitem{AK} M. Aoki, S. Koyama, \emph{Chebyshev's bias against splitting and principal primes in global fields}, J. Number Theory \textbf{245} (2023), 233--262. arXiv:2203.12266.

\bibitem{Booker} A. R. Booker, \emph{Artin's conjecture, Turing's method, and the Riemann hypothesis}, Experiment. Math. \textbf{15} (2006) (not read).

\bibitem{Cohn} J. H. E. Cohn, \emph{On square Fibonacci numbers}, J. London Math. Soc. \textbf{39} (1964), 537--540.

\bibitem{DJ} S. Dembner, V. Jain, \emph{Hyperelliptic curves and newform coefficients}, J. Number Theory (2021). arXiv:2007.08358.

\bibitem{Johnston} D. R. Johnston, arXiv:2201.06184 (2022); statement as quoted in \cite{Suzuki1} and in the arXiv abstract; paper not read.

\bibitem{Rankin} R. A. Rankin, \emph{Contributions to the theory of Ramanujan's function $\tau(n)$ and similar arithmetical functions. I}, Proc. Cambridge Philos. Soc. \textbf{35} (1939) (not consulted).

\bibitem{WR} R. Sharma, \emph{Rubinstein--Sarnak densities for weighted prime number races}, preprint, 2026, doi:10.5281/zenodo.21947101; and the working note \file{05-weighted-races-grh/explicit/} in \cite{Repo}.

\bibitem{Yoshida88} H. Yoshida, \emph{On calculations of zeros of $L$-functions related with Ramanujan's discriminant function on the critical line}, J. Ramanujan Math. Soc. \textbf{3} (1988), 87--95.

\bibitem{Angdinata} D. K. Angdinata et al.\ (24 authors), \emph{ABC implies that Ramanujan's tau function misses almost all primes}, arXiv:2603.29970 (v3, 2026).

\bibitem{BCO} J. S. Balakrishnan, W. Craig, K. Ono, \emph{Variations of Lehmer's conjecture for Ramanujan's tau-function}, J. Number Theory \textbf{237} (2022). arXiv:2005.10345.

\bibitem{BCOT} J. S. Balakrishnan, W. Craig, K. Ono, W.-L. Tsai, \emph{Variants of Lehmer's speculation for newforms}, Adv. Math. \textbf{428} (2023), Paper No.~109141; Corrigendum, Adv. Math. \textbf{435} (2023), 109347. arXiv:2005.10354 (v4 used).

\bibitem{BF} O. Balkanova, D. Frolenkov, \emph{The mean value of symmetric square $L$-functions}, Algebra Number Theory \textbf{12} (2018), 35--59.

\bibitem{BGPS} M. A. Bennett, A. Gherga, V. Patel, S. Siksek, \emph{Odd values of the Ramanujan tau function}, Math. Ann. \textbf{382} (2022). arXiv:2101.02933.

\bibitem{BHV} Yu. Bilu, G. Hanrot, P. M. Voutier, \emph{Existence of primitive divisors of Lucas and Lehmer numbers}, J. Reine Angew. Math. \textbf{539} (2001), 75--122.

\bibitem{BKLNW} S. Broadbent, H. Kadiri, A. Lumley, N. Ng, K. Wilk, \emph{Sharper bounds for the Chebyshev function $\theta(x)$}, Math. Comp. \textbf{90} (2021), 2281--2315. arXiv:2002.11068.

\bibitem{Deligne} P. Deligne, \emph{La conjecture de Weil. I}, Publ. Math. IH\'ES \textbf{43} (1974), 273--307.

\bibitem{Devin} L. Devin, \emph{Chebyshev's bias for analytic $L$-functions}, Math. Proc. Cambridge Philos. Soc. \textbf{169} (2020), 103--140. arXiv:1706.06394.

\bibitem{Dusart} P. Dusart, \emph{Estimates of some functions over primes without R.H.}, arXiv:1002.0442 (2010).

\bibitem{FKS} A. Fiori, H. Kadiri, J. Swidinsky, \emph{Sharper bounds for the Chebyshev function $\psi(x)$}, arXiv:2204.02588 (v3 read).

\bibitem{GJ} S. Gelbart, H. Jacquet, \emph{A relation between automorphic representations of $\mathrm{GL}(2)$ and $\mathrm{GL}(3)$}, Ann. Sci. \'Ec. Norm. Sup\'er. (4) \textbf{11} (1978), 471--542.

\bibitem{Goto} A. Goto, \emph{On some values which do not belong to the image of Ramanujan's tau-function}, arXiv:2405.16723v1 (2024); published in Arch. Math. \textbf{125} (2025), 157--172, doi:10.1007/s00013-025-02139-5 (published version not seen).

\bibitem{HM} K. Hanada, R. Madhukara, \emph{Fourier coefficients of level 1 Hecke eigenforms}, Acta Arith. \textbf{200} (2021). arXiv:2007.08683.

\bibitem{Hayani} M. Hayani, \emph{Chebyshev's bias without linear independence}, arXiv:2512.23302 (v3, 2026).

\bibitem{HIK} G. Hiary, S. Ireland, M. Kyi, \emph{A method for verifying the generalized Riemann hypothesis}, arXiv:2408.00187 (2024).

\bibitem{IK} H. Iwaniec, E. Kowalski, \emph{Analytic Number Theory}, Amer. Math. Soc. Colloq. Publ. \textbf{53}, AMS, Providence, 2004.

\bibitem{Arb} F. Johansson, \emph{Arb: efficient arbitrary-precision midpoint-radius interval arithmetic}, IEEE Trans. Comput. \textbf{66} (2017), 1281--1292.

\bibitem{KK} S. Koyama, N. Kurokawa, \emph{Chebyshev's bias for Ramanujan's $\tau$-function via the Deep Riemann Hypothesis}, Proc. Japan Acad. Ser. A Math. Sci. \textbf{98} (2022), no.~6, 35--39. arXiv:2203.12791.

\bibitem{KS} S. Koyama, A. Sheth, \emph{Chebyshev's bias for modular forms}, arXiv:2509.04187 (2025).

\bibitem{MV} H. L. Montgomery, R. C. Vaughan, \emph{Multiplicative Number Theory I. Classical Theory}, Cambridge Univ. Press, 2007.

\bibitem{PARI} The PARI Group, \emph{PARI/GP version 2.15.4}, Univ. Bordeaux, 2023, \url{http://pari.math.u-bordeaux.fr/}.

\bibitem{RS62} J. B. Rosser, L. Schoenfeld, \emph{Approximate formulas for some functions of prime numbers}, Illinois J. Math. \textbf{6} (1962), 64--94.

\bibitem{Sarnak} P. Sarnak, \emph{Letter to Barry Mazur on ``Chebyshev's bias'' for $\tau(p)$}, November 2007 (with an addendum of January 2008), \url{https://web.math.princeton.edu/sarnak/MazurLtrMay08.PDF}.

\bibitem{Sheth} A. Sheth, \emph{Euler products at the centre and applications to Chebyshev's bias}, Math. Proc. Cambridge Philos. Soc. \textbf{179} (2025), 331--349. arXiv:2405.01512.

\bibitem{Shimura} G. Shimura, \emph{On the holomorphy of certain Dirichlet series}, Proc. London Math. Soc. (3) \textbf{31} (1975), 79--98.

\bibitem{Suzuki1} M. Suzuki, \emph{On variants of Chebyshev's conjecture}, Ramanujan J. \textbf{68} (2025), Paper No.~95. arXiv:2411.07436.

\bibitem{Suzuki2} M. Suzuki, \emph{Screw functions of Dirichlet series in the extended Selberg class}, Int. J. Number Theory \textbf{21} (2025), 1815--1823. arXiv:2209.12832.

\bibitem{SwD} H. P. F. Swinnerton-Dyer, \emph{On $\ell$-adic representations and congruences for coefficients of modular forms}, in: Modular Functions of One Variable III, Lecture Notes in Math. \textbf{350}, Springer, 1973, 1--55.

\bibitem{Yoshida95} H. Yoshida, \emph{On calculations of zeros of various $L$-functions}, J. Math. Kyoto Univ. \textbf{35} (1995), no.~4. arXiv:math/9411213.

\bibitem{Zagier} D. Zagier, \emph{Modular forms whose Fourier coefficients involve zeta-functions of quadratic fields}, in: Modular Functions of One Variable VI, Lecture Notes in Math. \textbf{627}, Springer, 1977, 105--169.

\bibitem{Repo} R. Sharma, \emph{Research-Math} (code, certificates and working notes), \url{https://github.com/raghavs1729/Research-Math}, directory \repo.

\end{thebibliography}
\end{document}
